\documentclass[11pt]{article}

% ---------- Encodage / langue / polices ----------
\usepackage{fontspec}
\usepackage{polyglossia}
\setmainlanguage{french}
\usepackage{microtype}

% ---------- Mathématiques ----------
\usepackage{amsmath,amssymb,amsthm,mathtools,stmaryrd}
\usepackage{mathrsfs}

% ---------- Mise en page ----------
\usepackage[a4paper,margin=2.4cm]{geometry}
\usepackage{enumitem}
\usepackage{xcolor}
\usepackage{tikz}
\usetikzlibrary{arrows.meta}
\usepackage{booktabs}
\usepackage{array}
\usepackage[colorlinks=true,linkcolor=NavyBlue,citecolor=NavyBlue,urlcolor=NavyBlue]{hyperref}

% ---------- Couleurs ----------
\definecolor{NavyBlue}{RGB}{27,60,120}
\definecolor{ForestGreen}{RGB}{20,120,60}
\definecolor{Crimson}{RGB}{160,30,45}
\definecolor{BoxBlue}{RGB}{232,238,248}

% ---------- Styles de théorèmes ----------
\newtheoremstyle{thmblue}{10pt}{10pt}{\itshape}{}{\bfseries\color{NavyBlue}}{.}{.5em}{}
\newtheoremstyle{defplain}{10pt}{10pt}{}{}{\bfseries\color{ForestGreen}}{.}{.5em}{}
\newtheoremstyle{rem}{8pt}{8pt}{}{}{\itshape}{.}{.5em}{}

\theoremstyle{thmblue}
\newtheorem{theoreme}{Théorème}[section]
\newtheorem{proposition}[theoreme]{Proposition}
\newtheorem{lemme}[theoreme]{Lemme}
\newtheorem{corollaire}[theoreme]{Corollaire}

\theoremstyle{defplain}
\newtheorem{definition}[theoreme]{Définition}
\newtheorem{exemple}[theoreme]{Exemple}
\newtheorem{contreexemple}[theoreme]{Contre-exemple}

\theoremstyle{rem}
\newtheorem{remarque}[theoreme]{Remarque}
\newtheorem{notation}[theoreme]{Notation}

\renewcommand{\proofname}{\normalfont\bfseries D\'emonstration}

% ---------- Macros ----------
\newcommand{\R}{\mathbb{R}}
\newcommand{\N}{\mathbb{N}}
\newcommand{\Q}{\mathbb{Q}}
\newcommand{\dx}{\,dx}
\newcommand{\dt}{\,dt}
\newcommand{\du}{\,du}
\newcommand{\Var}{\operatorname{Var}}
\newcommand{\pp}{p.\,p.}
\DeclareMathOperator{\Lip}{Lip}
\DeclareMathOperator{\osc}{osc}
\DeclareMathOperator{\diam}{diam}
\newcommand{\abs}[1]{\left\lvert #1 \right\rvert}
\newcommand{\norm}[1]{\left\lVert #1 \right\rVert}
\newcommand{\eps}{\varepsilon}
\newcommand{\lstar}{\lambda^{\!*}}

\newcommand{\keybox}[2]{%
\begin{center}
\fcolorbox{NavyBlue}{BoxBlue}{\parbox{0.95\linewidth}{\textbf{#1}\\[2pt] #2}}
\end{center}}

\title{\vspace{-1.2cm}\bfseries\color{NavyBlue}\Large Le théorème fondamental de l'analyse\\[4pt]
\large de Riemann à Lebesgue}
\author{\normalsize Généralisations, contre-exemples et l'escalier de Cantor}
\date{\normalsize\today}

\begin{document}
\maketitle
\thispagestyle{empty}

\begin{abstract}
\noindent
On retrace la généralisation progressive du théorème fondamental de l'analyse (TFA), de sa forme
riemannienne classique jusqu'à sa formulation optimale dans le cadre de Lebesgue, puis au-delà. Le
fil conducteur est une série de \emph{contre-exemples} délimitant la frontière de validité de chaque
énoncé : la fonction de Volterra, puis \emph{l'escalier de Cantor}, enfin $x^2\sin(1/x^2)$. Toutes les
démonstrations sont données \emph{intégralement} et s'appuient sur une unique boîte à outils --- le
lemme de recouvrement de Vitali et les théorèmes de convergence de l'intégration de Lebesgue
(Fatou, convergence dominée) --- que l'on suppose connus. Aucun résultat n'est renvoyé à une annexe.
\end{abstract}

\tableofcontents
\vspace{0.4cm}

% ============================================================
\section{Introduction et boîte à outils}
% ============================================================

Sous sa forme scolaire, le TFA énonce, pour $f$ continue,
\[
\frac{d}{dx}\int_a^x f(t)\dt = f(x)
\qquad\text{et}\qquad
\int_a^b F'(t)\dt = F(b)-F(a)\quad(F'=f).
\]
Dès qu'on relâche la continuité, chacun de ces deux volets se fissure. La question centrale est :

\begin{center}
\emph{pour quelles fonctions $F$ a-t-on encore $\displaystyle\int_a^b F' = F(b)-F(a)$ ?}
\end{center}

La réponse, on le verra, est \emph{exactement} : les fonctions \emph{absolument continues}.
L'escalier de Cantor --- continu, croissant, de dérivée nulle presque partout, mais croissant de $0$
à $1$ --- est le témoin qui interdit toute réponse plus laxiste.

\paragraph{Conventions.} On travaille sur un segment $[a,b]$. La lettre $\lambda$ désigne la mesure de
Lebesgue et $\lstar$ la mesure extérieure ; « \pp » signifie « pour $\lambda$-presque tout $x$ ».
On note $L^1=L^1([a,b])$ l'espace des fonctions Lebesgue-intégrables. On suppose acquis, du côté de
l'\emph{intégration} : la construction de l'intégrale de Lebesgue et de la tribu des ensembles
mesurables, le lemme de Fatou, les théorèmes de convergence monotone et dominée. Du côté de la
\emph{mesure}, on ne suppose que la définition de $\lambda$ sur les mesurables ; les propriétés de la
mesure extérieure $\lstar$ dont on se sert (sous-additivité, régularité extérieure, continuité
croissante) sont rappelées et démontrées au §\ref{sub:mesext} ci-dessous.

\medskip
Les démonstrations « fines » reposent sur deux ingrédients géométriques : les propriétés de la mesure
extérieure, puis un lemme de recouvrement. On les établit dans cet ordre.

\subsection{Rappel : la mesure extérieure de Lebesgue}\label{sub:mesext}

\begin{definition}[mesure extérieure]
Pour $E\subset\R$, la \emph{mesure extérieure de Lebesgue} est
\[
\lstar(E)=\inf\Bigl\{\sum_{n\ge1}\ell(I_n)\ :\ E\subset\bigcup_{n\ge1}I_n,\ I_n\text{ intervalles ouverts}\Bigr\},
\]
où $\ell(]c,d[)=d-c$. C'est un infimum sur tous les recouvrements dénombrables de $E$ par des
intervalles ouverts.
\end{definition}

De la définition découlent immédiatement la \emph{positivité} ($\lstar\ge0$, $\lstar(\varnothing)=0$),
la \emph{monotonie} ($E\subset E'\Rightarrow\lstar(E)\le\lstar(E')$, tout recouvrement de $E'$ en étant
un de $E$) et l'\emph{invariance par translation}. On admet le fait de base
$\lstar([c,d])=\lstar(]c,d[)=d-c$ (l'inégalité $\le$ est immédiate ; l'inégalité $\ge$ résulte, pour un
recouvrement d'un segment, de la compacité de Borel--Lebesgue). Trois propriétés servent constamment.

\begin{proposition}[sous-additivité dénombrable]\label{prop:sousadd}
Pour toute suite $(A_n)$ de parties de $\R$,
\[
\lstar\Bigl(\bigcup_n A_n\Bigr)\le\sum_n\lstar(A_n).
\]
\end{proposition}

\begin{proof}
Si un $\lstar(A_n)=\infty$, le membre de droite est infini et il n'y a rien à montrer. Sinon, fixons
$\eps>0$. Par définition de l'infimum, pour chaque $n$ il existe des intervalles ouverts
$(I_{n,j})_j$ recouvrant $A_n$ avec $\sum_j\ell(I_{n,j})<\lstar(A_n)+\eps\,2^{-n}$. La famille
dénombrable $(I_{n,j})_{n,j}$ recouvre $\bigcup_n A_n$, d'où
\[
\lstar\Bigl(\bigcup_n A_n\Bigr)\le\sum_{n,j}\ell(I_{n,j})<\sum_n\bigl(\lstar(A_n)+\eps\,2^{-n}\bigr)
=\sum_n\lstar(A_n)+\eps.
\]
On fait $\eps\to0$.
\end{proof}

\begin{proposition}[régularité extérieure]\label{prop:regext}
Pour tout $E\subset\R$ et tout $\eps>0$, il existe un ouvert $O\supset E$ tel que
$\lambda(O)\le\lstar(E)+\eps$. En particulier $\lstar(E)=\inf\{\lambda(O):O\text{ ouvert},\ O\supset E\}$.
\end{proposition}

\begin{proof}
Si $\lstar(E)=\infty$, prendre $O=\R$. Sinon, par définition de l'infimum, il existe des intervalles
ouverts $(I_n)$ recouvrant $E$ avec $\sum_n\ell(I_n)<\lstar(E)+\eps$. L'ensemble $O=\bigcup_n I_n$ est
ouvert, contient $E$, et par sous-additivité (proposition~\ref{prop:sousadd})
\[
\lambda(O)=\lstar(O)\le\sum_n\lstar(I_n)=\sum_n\ell(I_n)<\lstar(E)+\eps.\qedhere
\]
\end{proof}

\begin{corollaire}[enveloppe mesurable]\label{cor:enveloppe}
Tout $A\subset\R$ possède une \emph{enveloppe mesurable} : un ensemble mesurable (même de type
$G_\delta$) $H\supset A$ tel que $\lambda(H)=\lstar(A)$.
\end{corollaire}

\begin{proof}
Par la proposition~\ref{prop:regext}, choisissons pour chaque $n\ge1$ un ouvert $O_n\supset A$ avec
$\lambda(O_n)\le\lstar(A)+\tfrac1n$. Alors $H=\bigcap_n O_n$ est un $G_\delta$ (donc mesurable),
$A\subset H$, et pour tout $n$ : $\lstar(A)\le\lambda(H)\le\lambda(O_n)\le\lstar(A)+\tfrac1n$. En
faisant $n\to\infty$, $\lambda(H)=\lstar(A)$.
\end{proof}

\begin{proposition}[continuité croissante de $\lstar$]\label{prop:contcroiss}
Si $E_1\subset E_2\subset\cdots$ et $E=\bigcup_n E_n$, alors $\lstar(E)=\lim_n\lstar(E_n)$.
\end{proposition}

\begin{proof}
La suite $\bigl(\lstar(E_n)\bigr)$ est croissante (monotonie), et $\lstar(E)\ge\lstar(E_n)$ pour tout
$n$, donc $\lstar(E)\ge\lim_n\lstar(E_n)$. Pour l'inégalité inverse, prenons pour chaque $n$ une
enveloppe mesurable $H_n\supset E_n$ avec $\lambda(H_n)=\lstar(E_n)$
(corollaire~\ref{cor:enveloppe}), puis posons $G_n=\bigcap_{k\ge n}H_k$. Chaque $G_n$ est mesurable ;
comme $E_n\subset E_k\subset H_k$ pour $k\ge n$, on a $E_n\subset G_n$, d'où
$\lstar(E_n)\le\lambda(G_n)\le\lambda(H_n)=\lstar(E_n)$, soit $\lambda(G_n)=\lstar(E_n)$. La suite
$(G_n)$ est croissante et $G:=\bigcup_n G_n\supset\bigcup_n E_n=E$. Par continuité croissante de la
mesure $\lambda$ sur les ensembles mesurables (additivité dénombrable),
\[
\lstar(E)\le\lambda(G)=\lim_n\lambda(G_n)=\lim_n\lstar(E_n).\qedhere
\]
\end{proof}

\begin{remarque}[ensembles mesurables]
Rappelons qu'un ensemble $E$ est \emph{Lebesgue-mesurable} (au sens de Carathéodory) si
$\lstar(T)=\lstar(T\cap E)+\lstar(T\setminus E)$ pour tout $T\subset\R$. Les mesurables forment une
tribu contenant les ouverts (donc les boréliens), sur laquelle $\lstar$ est \emph{dénombrablement
additive} ; sa restriction est la mesure de Lebesgue $\lambda$, complète. Sur un mesurable,
$\lambda=\lstar$ : c'est pourquoi on écrit $\lambda$ dès que l'argument porte sur un ouvert ou un
ensemble mesurable, et $\lstar$ quand on manipule un ensemble a priori quelconque (typiquement les
ensembles $E_{s,r}$ de la preuve du théorème~\ref{thm:mono}, dont on montre justement qu'ils sont
négligeables).
\end{remarque}

\keybox{Mode d'emploi : $\lambda$ ou $\lstar$ ?}{%
Devant un ensemble, une seule question : \emph{le sait-on mesurable ?}
\\[3pt]
$\bullet$~\textbf{Oui} (intervalle, ouvert, fermé, borélien, ou construit à partir d'eux par
opérations dénombrables) : on écrit $\lambda$, et l'on a le droit de \emph{découper et sommer}
(additivité dénombrable). C'est le régime des ouverts $O,U$ et des intervalles $I_k$, là où l'on veut
des \emph{égalités} de longueurs comme $\lambda(O)=\sum_n\ell(I_n)$.
\\[3pt]
$\bullet$~\textbf{Non} / ensemble a priori quelconque / \emph{image directe} $H(A)$ / ensemble que l'on
cherche justement à déclarer négligeable (les $E_{s,r}$) : seul $\lstar$ est défini, avec les seules
règles \emph{faibles} --- monotonie, sous-additivité, régularité, continuité croissante --- donc au
mieux des \emph{inégalités}.
\\[5pt]
\textbf{Le pont (complétude).} Dès qu'on a établi $\lstar(E)=0$, l'ensemble \emph{devient} mesurable
et $\lambda$ prend le relais :
\[
\lstar(E)=0\ \Longrightarrow\ E\ \text{mesurable}\ \Longrightarrow\ \lambda(E)=0.
\]
Chaque démonstration « fine » suit ce mouvement : partir de $\lstar$ sur un ensemble suspect, montrer
qu'il est nul, puis basculer sans douleur dans le monde de $\lambda$ --- indolore car $\lambda=\lstar$
sur les mesurables.}

\subsection{Le lemme de recouvrement de Vitali}

\begin{definition}[Recouvrement de Vitali]
Soit $E\subset\R$. Une famille $\mathcal V$ d'intervalles fermés de longueur $>0$ \emph{recouvre $E$
au sens de Vitali} si, pour tout $x\in E$ et tout $\rho>0$, il existe $I\in\mathcal V$ tel que
$x\in I$ et $0<\lambda(I)<\rho$.
\end{definition}

\begin{lemme}[Vitali]\label{lem:vitali}
Soit $E\subset\R$ avec $\lstar(E)<\infty$, et $\mathcal V$ un recouvrement de $E$ au sens de Vitali.
Pour tout $\eps>0$, il existe une famille \emph{finie} d'intervalles \emph{deux à deux disjoints}
$I_1,\dots,I_N\in\mathcal V$ tels que
\[
\lstar\!\Bigl(E\setminus\bigcup_{k=1}^{N} I_k\Bigr)<\eps.
\]
\end{lemme}

\begin{proof}
Par régularité extérieure (proposition~\ref{prop:regext}), fixons un ouvert $O\supset E$ avec $\lambda(O)<\infty$. Quitte à ne garder
que les $I\in\mathcal V$ inclus dans $O$ (ce qui reste un recouvrement de Vitali de $E$, car les
intervalles peuvent être pris arbitrairement petits autour de chaque point de $E\subset O$ ouvert),
on suppose $I\subset O$ pour tout $I\in\mathcal V$.

On construit une suite disjointe par sélection gloutonne. Choisissons $I_1\in\mathcal V$ quelconque.
Supposons $I_1,\dots,I_n$ choisis, deux à deux disjoints. Si $E\subset\bigcup_{k\le n}I_k$, on
s'arrête et le résultat est acquis. Sinon, posons
\[
s_n=\sup\{\lambda(I):I\in\mathcal V,\ I\cap I_k=\varnothing\ \forall k\le n\}.
\]
L'ensemble est non vide (un point $x\in E\setminus\bigcup_{k\le n}I_k$, qui est ouvert, admet de
petits $I\in\mathcal V$ disjoints des $I_k$), et $s_n\le\lambda(O)<\infty$. Choisissons
$I_{n+1}\in\mathcal V$ disjoint des précédents avec $\lambda(I_{n+1})>s_n/2$.

Si le procédé ne s'arrête pas, on obtient une suite infinie $(I_k)$ deux à deux disjointes incluses
dans $O$, donc $\sum_k\lambda(I_k)\le\lambda(O)<\infty$ ; en particulier $\lambda(I_k)\to0$ et le
reste $\sum_{k>N}\lambda(I_k)\to0$.

\emph{Estimation du reste.} Montrons : pour tout $N$,
$\displaystyle E\setminus\bigcup_{k\le N}I_k\subset\bigcup_{k>N}5I_k$, où $5I_k$ désigne l'intervalle
de même centre que $I_k$ et de longueur $5\lambda(I_k)$. Soit
$x\in E\setminus\bigcup_{k\le N}I_k$. Comme $\bigcup_{k\le N}I_k$ est fermé et ne contient pas $x$,
il existe $I\in\mathcal V$ avec $x\in I$, disjoint de $I_1,\dots,I_N$. Cet $I$ doit rencontrer l'un
des $I_k$ ($k>N$) : sinon $I$ serait éligible à toutes les étapes, d'où
$\lambda(I)\le s_n<2\lambda(I_{n+1})\to0$, contredisant $\lambda(I)>0$ fixe. Soit $k$ le plus petit
indice tel que $I\cap I_k\neq\varnothing$ ; alors $k>N$ et $I$ est disjoint de $I_1,\dots,I_{k-1}$,
donc $\lambda(I)\le s_{k-1}<2\lambda(I_k)$. Comme $I$ rencontre $I_k$ et que
$\lambda(I)<2\lambda(I_k)$, tout point de $I$ est à distance $<\lambda(I)<2\lambda(I_k)$ d'un point de
$I_k$ : donc $I\subset 5I_k$, et en particulier $x\in5I_k$.

Ainsi $\lstar\bigl(E\setminus\bigcup_{k\le N}I_k\bigr)\le\sum_{k>N}5\lambda(I_k)\xrightarrow[N\to\infty]{}0$.
On choisit $N$ rendant ce reste $<\eps$.
\end{proof}

Ce lemme, et les théorèmes de convergence de Lebesgue, suffisent à tout démontrer dans ce texte.

\keybox{Feuille de route}{%
Mesure extérieure $\to$ Vitali $\to$ dérivabilité \pp{} des fonctions monotones $\to$ escalier de Cantor $\to$ variation
bornée \& continuité absolue $\to$ lemme d'image \& lemme de rigidité $\to$ les deux volets du TFA de
Lebesgue $\to$ décomposition d'une fonction $BV$ $\to$ intégrale de Henstock--Kurzweil.}

% ============================================================
\section{Le théorème fondamental classique, et sa première fissure}
% ============================================================

\begin{theoreme}[TFA classique]\label{thm:classique}
Soit $f:[a,b]\to\R$ continue.
\begin{enumerate}[label=\textup{(\alph*)},leftmargin=2.2em,itemsep=1pt]
  \item La fonction $F(x)=\int_a^x f(t)\dt$ est de classe $C^1$, et $F'=f$.
  \item Si $G'=f$ partout, alors $\int_a^b f(t)\dt = G(b)-G(a)$.
\end{enumerate}
\end{theoreme}

\begin{proof}
(a) Soit $x_0\in[a,b]$ et $\eps>0$. La continuité de $f$ en $x_0$ fournit $\delta>0$ tel que
$\abs{f(t)-f(x_0)}<\eps$ pour $\abs{t-x_0}<\delta$. Pour $0<\abs h<\delta$,
\[
\abs*{\frac{F(x_0+h)-F(x_0)}{h}-f(x_0)}
=\abs*{\frac1h\int_{x_0}^{x_0+h}\!\bigl(f(t)-f(x_0)\bigr)\dt}
\le\frac1{\abs h}\,\abs h\,\eps=\eps,
\]
donc $F'(x_0)=f(x_0)$ ; $F'=f$ est continue, d'où $F\in C^1$. (b) $G-F$ a une dérivée nulle sur
l'intervalle $[a,b]$, donc est constante (théorème des accroissements finis) ; en évaluant en $a$ et
$b$ on obtient $G(b)-G(a)=F(b)-F(a)=\int_a^b f$.
\end{proof}

Peut-on affaiblir (b) en supposant seulement $G$ \emph{dérivable partout, de dérivée bornée} ? Dans le
cadre de Riemann, non.

\begin{contreexemple}[Volterra, 1881]\label{ce:volterra}
Il existe $V:[0,1]\to\R$ dérivable en tout point, à dérivée $V'$ bornée, telle que $V'$ n'est
\emph{pas Riemann-intégrable}. On part d'un ensemble de Cantor \emph{gras} $K^\sharp$, de mesure
$\lambda(K^\sharp)>0$, obtenu en retirant à chaque étape des intervalles centraux de longueurs
sommables $<1$. Sur chaque intervalle contigu $]a_j,b_j[$ on greffe une copie tronquée de
$g(x)=x^2\sin(1/x)$ (dérivable, dérivée bornée, oscillant infiniment près des bords), recollée pour
que $V$ soit dérivable jusqu'aux extrémités avec $V'=0$ sur $K^\sharp$. En chaque point de $K^\sharp$,
$V'$ oscille et est discontinue ; l'ensemble des discontinuités de $V'$ contient $K^\sharp$, donc est
de mesure $>0$. Le critère de Lebesgue (une fonction bornée est Riemann-intégrable si et seulement si
l'ensemble de ses discontinuités est négligeable) montre que $V'$ n'est pas Riemann-intégrable.
\end{contreexemple}

\keybox{Diagnostic}{%
L'intégrale de Riemann est trop étroite : l'énoncé « $G$ dérivable $\Rightarrow\int G'=G(b)-G(a)$ »
n'a pas toujours de sens, faute d'intégrabilité de $G'$. Il faut une intégrale plus robuste --- celle
de Lebesgue --- dont le premier acquis est un puissant théorème de différentiation.}

% ============================================================
\section{Dérivabilité presque partout des fonctions monotones}
% ============================================================

Pour une fonction $f$ quelconque, on introduit les quatre \emph{dérivées de Dini}
\[
D^{+}f(x)=\limsup_{h\to0^+}\frac{f(x+h)-f(x)}{h},\quad
D_{+}f(x)=\liminf_{h\to0^+}\frac{f(x+h)-f(x)}{h},
\]
et de même $D^{-}f,D_{-}f$ pour $h\to0^-$. La fonction $f$ est dérivable en $x$ (au sens fini) si et
seulement si ces quatre nombres coïncident et sont finis.

\begin{theoreme}[Lebesgue, 1904]\label{thm:mono}
Soit $f:[a,b]\to\R$ croissante. Alors $f$ est dérivable \pp, sa dérivée $f'$ est mesurable, positive,
et
\[
\int_a^b f'(x)\dx\ \le\ f(b)-f(a).
\]
\end{theoreme}

\begin{proof}
\emph{Étape 1 : $f$ est dérivable \pp.} Il suffit de montrer que, presque partout, les quatre dérivées
de Dini coïncident et $D^{+}f<\infty$. Par un argument symétrique (appliqué à $x\mapsto-f(-x)$), il
suffit de traiter l'inégalité $D^{+}f\le D_{-}f$ \pp{} (l'inégalité inverse $D_{-}f\le D^{+}f$ étant
toujours plus facile ; en la combinant à ses analogues on obtient l'égalité des quatre). Fixons deux
rationnels $0<s<r$ et posons
\[
E_{s,r}=\{x\in(a,b):D_{-}f(x)<s<r<D^{+}f(x)\}.
\]
L'ensemble où $D^{+}f>D_{-}f$ est la réunion \emph{dénombrable} $\bigcup_{s<r}E_{s,r}$ ; il suffit
donc de prouver $\lstar(E_{s,r})=0$. Notons $\alpha=\lstar(E_{s,r})$ et fixons $\eps>0$. Par
régularité extérieure (proposition~\ref{prop:regext}), soit $O\supset E_{s,r}$ ouvert avec
$\lambda(O)<\alpha+\eps$.

Pour $x\in E_{s,r}$, l'inégalité $D_{-}f(x)<s$ fournit des $h>0$ arbitrairement petits tels que
$[x-h,x]\subset O$ et $f(x)-f(x-h)<s\,h$. Ces intervalles forment un recouvrement de Vitali de
$E_{s,r}$ ; par le lemme~\ref{lem:vitali}, on en extrait une famille finie disjointe
$[x_i-h_i,x_i]$, $i=1,\dots,p$, telle que $\lstar\bigl(E_{s,r}\cap\bigcup_i(x_i-h_i,x_i)\bigr)>\alpha-\eps$.
En sommant,
\begin{equation}\label{eq:mono1}
\sum_{i=1}^{p}\bigl(f(x_i)-f(x_i-h_i)\bigr)<s\sum_{i=1}^{p}h_i\le s\,\lambda(O)<s(\alpha+\eps).
\end{equation}
Posons $A=E_{s,r}\cap\bigcup_i(x_i-h_i,x_i)$, de sorte que $\lstar(A)>\alpha-\eps$. Pour $y\in A$,
l'inégalité $D^{+}f(y)>r$ fournit des $k>0$ arbitrairement petits tels que $[y,y+k]$ soit inclus dans
l'un des intervalles ouverts $(x_i-h_i,x_i)$ et $f(y+k)-f(y)>r\,k$. Nouveau recouvrement de Vitali de
$A$, nouvelle extraction finie disjointe $[y_j,y_j+k_j]$, $j=1,\dots,q$, avec
$\sum_j k_j>\lstar(A)-\eps>\alpha-2\eps$. En sommant,
\begin{equation}\label{eq:mono2}
\sum_{j=1}^{q}\bigl(f(y_j+k_j)-f(y_j)\bigr)>r\sum_{j=1}^{q}k_j>r(\alpha-2\eps).
\end{equation}
Chaque $[y_j,y_j+k_j]$ est contenu dans un $(x_i-h_i,x_i)$ ; comme $f$ est croissante, la somme des
accroissements de $f$ sur les $[y_j,y_j+k_j]$ contenus dans un même $(x_i-h_i,x_i)$ est majorée par
l'accroissement $f(x_i)-f(x_i-h_i)$. En regroupant,
\[
\sum_{j=1}^{q}\bigl(f(y_j+k_j)-f(y_j)\bigr)\ \le\ \sum_{i=1}^{p}\bigl(f(x_i)-f(x_i-h_i)\bigr).
\]
Combinée à \eqref{eq:mono1} et \eqref{eq:mono2}, cela donne $r(\alpha-2\eps)<s(\alpha+\eps)$. En
faisant $\eps\to0$ : $r\alpha\le s\alpha$, soit $(r-s)\alpha\le0$. Comme $r>s$, on a $\alpha=0$.
Ainsi $D^{+}f=D_{-}f$ \pp, et $f$ est dérivable \pp.

\emph{Étape 2 : intégrabilité et inégalité.} Prolongeons $f$ par $f(x)=f(b)$ pour $x>b$ et posons
$g_n(x)=n\bigl(f(x+\tfrac1n)-f(x)\bigr)\ge0$. Les $g_n$ sont mesurables (comme $f$ croissante l'est)
et $g_n\to f'$ \pp. Par le lemme de Fatou,
\[
\int_a^b f'\ \le\ \liminf_n\int_a^b g_n
= \liminf_n\ n\!\int_a^b\!\bigl(f(x+\tfrac1n)-f(x)\bigr)\dx
= \liminf_n\Bigl[n\!\int_b^{b+\frac1n}\!\!f - n\!\int_a^{a+\frac1n}\!\!f\Bigr].
\]
Or $n\int_b^{b+1/n}f=f(b)$ (car $f\equiv f(b)$ au-delà de $b$) et
$n\int_a^{a+1/n}f\ge f(a)$ (car $f\ge f(a)$ sur $[a,a+1/n]$). Donc
$\int_a^b f'\le f(b)-f(a)<\infty$. En particulier $f'<\infty$ \pp.
\end{proof}

\begin{remarque}[le point crucial]\label{rem:strict}
L'inégalité peut être \emph{strictement} inférieure. Une fonction croissante peut « gagner de la
hauteur » sans que sa dérivée en rende compte. Le cas d'égalité est exactement ce que le TFA de
Lebesgue caractérisera ; l'escalier de Cantor réalise l'extrême opposé : $\int_0^1 f'=0$ pour
$f(1)-f(0)=1$.
\end{remarque}

% ============================================================
\section{L'escalier de Cantor}
% ============================================================

\subsection{L'ensemble triadique et sa mesure}

\begin{definition}
$K_0=[0,1]$ ; $K_{n+1}$ s'obtient en retirant le tiers ouvert central de chacun des $2^n$ intervalles
de $K_n$. L'\emph{ensemble de Cantor} est $K=\bigcap_{n\ge0}K_n$.
\end{definition}

\begin{center}
\begin{tikzpicture}[y=0pt,x=12cm]
  \foreach \lvl/\segs in {%
     0/{0/1},
     1/{0/0.3333, 0.6667/1},
     2/{0/0.1111,0.2222/0.3333,0.6667/0.7778,0.8889/1},
     3/{0/0.0370,0.0741/0.1111,0.2222/0.2593,0.2963/0.3333,0.6667/0.7037,0.7407/0.7778,0.8889/0.9259,0.9630/1},
     4/{0/0.0123,0.0247/0.0370,0.0741/0.0864,0.0988/0.1111,0.2222/0.2346,0.2469/0.2593,0.2963/0.3086,0.3210/0.3333,0.6667/0.6790,0.6914/0.7037,0.7407/0.7531,0.7654/0.7778,0.8889/0.9012,0.9136/0.9259,0.9630/0.9753,0.9877/1}%
  }{
    \foreach \s/\e in \segs {\draw[line width=3pt, NavyBlue] (\s,{-15*\lvl}) -- (\e,{-15*\lvl});}
    \node[left] at (-0.02,{-15*\lvl}) {\scriptsize $K_{\lvl}$};
  }
\end{tikzpicture}
\end{center}

\noindent
$K$ est compact, d'intérieur vide (il ne contient aucun intervalle), non dénombrable (les points de
$K$ sont les réels admettant un développement triadique n'utilisant que les chiffres $0$ et $2$, en
bijection avec $\{0,2\}^{\N}$). La longueur retirée vaut
$\sum_{n\ge0}2^n\cdot3^{-(n+1)}=\tfrac13\sum_{n\ge0}(2/3)^n=1$, donc $\lambda(K)=0$. Le complémentaire
$[0,1]\setminus K$ est une réunion dénombrable d'intervalles ouverts disjoints, de mesure totale $1$.

\subsection{Construction et propriétés de la fonction}

\begin{definition}\label{def:cantor}
Pour $x=\sum_{i\ge1}a_i3^{-i}$ ($a_i\in\{0,1,2\}$), soit $N$ le plus petit indice avec $a_N=1$
($N=\infty$ sinon). On pose $b_i=a_i/2$ pour $i<N$, $b_N=1$, $b_i=0$ pour $i>N$, et
\[
c(x)=\sum_{i\ge1}b_i\,2^{-i}.
\]
De manière équivalente : $c$ est \emph{constante} sur chaque intervalle contigu de $K$, égale à la
valeur commune à ses extrémités ; sur $K$, on remplace chaque chiffre triadique $0$ ou $2$ par le
chiffre binaire $0$ ou $1$.
\end{definition}

Concrètement $c=\tfrac12$ sur $[\tfrac13,\tfrac23]$, $c=\tfrac14$ sur $[\tfrac19,\tfrac29]$,
$c=\tfrac34$ sur $[\tfrac79,\tfrac89]$, etc.

\begin{figure}[h]
\centering
\begin{tikzpicture}[x=8cm,y=5cm]
  \draw[->,gray!70] (-0.02,0) -- (1.06,0) node[right]{\scriptsize $x$};
  \draw[->,gray!70] (0,-0.02) -- (0,1.08) node[above]{\scriptsize $c(x)$};
  \foreach \t in {0,0.3333,0.6667,1} \draw[gray!40] (\t,0)--(\t,-0.012) node[below]{\scriptsize $\t$};
  \foreach \h/\lab in {0/0,0.5/{1/2},1/1} \draw[gray!40] (0,\h)--(-0.012,\h) node[left]{\scriptsize $\lab$};
  \draw[dashed,gray!50] (0,1)--(1,1);
  \draw[line width=0.9pt, ForestGreen] plot coordinates {
(0.0000,0.0000) (0.0007,0.0078) (0.0013,0.0143) (0.0020,0.0156) (0.0027,0.0156) (0.0033,0.0234)
(0.0040,0.0286) (0.0047,0.0312) (0.0053,0.0312) (0.0060,0.0312) (0.0067,0.0312) (0.0073,0.0312)
(0.0080,0.0312) (0.0087,0.0381) (0.0093,0.0430) (0.0100,0.0469) (0.0107,0.0469) (0.0113,0.0527)
(0.0120,0.0571) (0.0127,0.0625) (0.0133,0.0625) (0.0140,0.0625) (0.0147,0.0625) (0.0153,0.0625)
(0.0160,0.0625) (0.0167,0.0625) (0.0173,0.0625) (0.0180,0.0625) (0.0187,0.0625) (0.0193,0.0625)
(0.0200,0.0625) (0.0207,0.0625) (0.0213,0.0625) (0.0220,0.0625) (0.0227,0.0625) (0.0233,0.0625)
(0.0240,0.0625) (0.0247,0.0625) (0.0253,0.0703) (0.0260,0.0762) (0.0267,0.0781) (0.0273,0.0781)
(0.0280,0.0859) (0.0287,0.0905) (0.0293,0.0938) (0.0300,0.0938) (0.0307,0.0938) (0.0313,0.0938)
(0.0320,0.0938) (0.0327,0.0938) (0.0333,0.1000) (0.0340,0.1055) (0.0347,0.1094) (0.0353,0.1094)
(0.0360,0.1143) (0.0367,0.1191) (0.0373,0.1250) (0.0380,0.1250) (0.0387,0.1250) (0.0393,0.1250)
(0.0400,0.1250) (0.0407,0.1250) (0.0413,0.1250) (0.0420,0.1250) (0.0427,0.1250) (0.0433,0.1250)
(0.0440,0.1250) (0.0447,0.1250) (0.0453,0.1250) (0.0460,0.1250) (0.0467,0.1250) (0.0473,0.1250)
(0.0480,0.1250) (0.0487,0.1250) (0.0493,0.1250) (0.0500,0.1250) (0.0507,0.1250) (0.0513,0.1250)
(0.0520,0.1250) (0.0527,0.1250) (0.0533,0.1250) (0.0540,0.1250) (0.0547,0.1250) (0.0553,0.1250)
(0.0560,0.1250) (0.0567,0.1250) (0.0573,0.1250) (0.0580,0.1250) (0.0587,0.1250) (0.0593,0.1250)
(0.0600,0.1250) (0.0607,0.1250) (0.0613,0.1250) (0.0620,0.1250) (0.0627,0.1250) (0.0633,0.1250)
(0.0640,0.1250) (0.0647,0.1250) (0.0653,0.1250) (0.0660,0.1250) (0.0667,0.1250) (0.0673,0.1250)
(0.0680,0.1250) (0.0687,0.1250) (0.0693,0.1250) (0.0700,0.1250) (0.0707,0.1250) (0.0713,0.1250)
(0.0720,0.1250) (0.0727,0.1250) (0.0733,0.1250) (0.0740,0.1250) (0.0747,0.1328) (0.0753,0.1382)
(0.0760,0.1406) (0.0767,0.1406) (0.0773,0.1484) (0.0780,0.1523) (0.0787,0.1562) (0.0793,0.1562)
(0.0800,0.1562) (0.0807,0.1562) (0.0813,0.1562) (0.0820,0.1562) (0.0827,0.1621) (0.0833,0.1667)
(0.0840,0.1719) (0.0847,0.1719) (0.0853,0.1758) (0.0860,0.1809) (0.0867,0.1875) (0.0873,0.1875)
(0.0880,0.1875) (0.0887,0.1875) (0.0893,0.1875) (0.0900,0.1875) (0.0907,0.1875) (0.0913,0.1875)
(0.0920,0.1875) (0.0927,0.1875) (0.0933,0.1875) (0.0940,0.1875) (0.0947,0.1875) (0.0953,0.1875)
(0.0960,0.1875) (0.0967,0.1875) (0.0973,0.1875) (0.0980,0.1875) (0.0987,0.1875) (0.0993,0.1953)
(0.1000,0.2000) (0.1007,0.2031) (0.1013,0.2031) (0.1020,0.2109) (0.1027,0.2148) (0.1033,0.2188)
(0.1040,0.2188) (0.1047,0.2188) (0.1053,0.2188) (0.1060,0.2188) (0.1067,0.2188) (0.1073,0.2236)
(0.1080,0.2285) (0.1087,0.2344) (0.1093,0.2344) (0.1100,0.2383) (0.1107,0.2429) (0.1113,0.2500)
(0.1120,0.2500) (0.1127,0.2500) (0.1133,0.2500) (0.1140,0.2500) (0.1147,0.2500) (0.1153,0.2500)
(0.1160,0.2500) (0.1167,0.2500) (0.1173,0.2500) (0.1180,0.2500) (0.1187,0.2500) (0.1193,0.2500)
(0.1200,0.2500) (0.1207,0.2500) (0.1213,0.2500) (0.1220,0.2500) (0.1227,0.2500) (0.1233,0.2500)
(0.1240,0.2500) (0.1247,0.2500) (0.1253,0.2500) (0.1260,0.2500) (0.1267,0.2500) (0.1273,0.2500)
(0.1280,0.2500) (0.1287,0.2500) (0.1293,0.2500) (0.1300,0.2500) (0.1307,0.2500) (0.1313,0.2500)
(0.1320,0.2500) (0.1327,0.2500) (0.1333,0.2500) (0.1340,0.2500) (0.1347,0.2500) (0.1353,0.2500)
(0.1360,0.2500) (0.1367,0.2500) (0.1373,0.2500) (0.1380,0.2500) (0.1387,0.2500) (0.1393,0.2500)
(0.1400,0.2500) (0.1407,0.2500) (0.1413,0.2500) (0.1420,0.2500) (0.1427,0.2500) (0.1433,0.2500)
(0.1440,0.2500) (0.1447,0.2500) (0.1453,0.2500) (0.1460,0.2500) (0.1467,0.2500) (0.1473,0.2500)
(0.1480,0.2500) (0.1487,0.2500) (0.1493,0.2500) (0.1500,0.2500) (0.1507,0.2500) (0.1513,0.2500)
(0.1520,0.2500) (0.1527,0.2500) (0.1533,0.2500) (0.1540,0.2500) (0.1547,0.2500) (0.1553,0.2500)
(0.1560,0.2500) (0.1567,0.2500) (0.1573,0.2500) (0.1580,0.2500) (0.1587,0.2500) (0.1593,0.2500)
(0.1600,0.2500) (0.1607,0.2500) (0.1613,0.2500) (0.1620,0.2500) (0.1627,0.2500) (0.1633,0.2500)
(0.1640,0.2500) (0.1647,0.2500) (0.1653,0.2500) (0.1660,0.2500) (0.1667,0.2500) (0.1673,0.2500)
(0.1680,0.2500) (0.1687,0.2500) (0.1693,0.2500) (0.1700,0.2500) (0.1707,0.2500) (0.1713,0.2500)
(0.1720,0.2500) (0.1727,0.2500) (0.1733,0.2500) (0.1740,0.2500) (0.1747,0.2500) (0.1753,0.2500)
(0.1760,0.2500) (0.1767,0.2500) (0.1773,0.2500) (0.1780,0.2500) (0.1787,0.2500) (0.1793,0.2500)
(0.1800,0.2500) (0.1807,0.2500) (0.1813,0.2500) (0.1820,0.2500) (0.1827,0.2500) (0.1833,0.2500)
(0.1840,0.2500) (0.1847,0.2500) (0.1853,0.2500) (0.1860,0.2500) (0.1867,0.2500) (0.1873,0.2500)
(0.1880,0.2500) (0.1887,0.2500) (0.1893,0.2500) (0.1900,0.2500) (0.1907,0.2500) (0.1913,0.2500)
(0.1920,0.2500) (0.1927,0.2500) (0.1933,0.2500) (0.1940,0.2500) (0.1947,0.2500) (0.1953,0.2500)
(0.1960,0.2500) (0.1967,0.2500) (0.1973,0.2500) (0.1980,0.2500) (0.1987,0.2500) (0.1993,0.2500)
(0.2000,0.2500) (0.2007,0.2500) (0.2013,0.2500) (0.2020,0.2500) (0.2027,0.2500) (0.2033,0.2500)
(0.2040,0.2500) (0.2047,0.2500) (0.2053,0.2500) (0.2060,0.2500) (0.2067,0.2500) (0.2073,0.2500)
(0.2080,0.2500) (0.2087,0.2500) (0.2093,0.2500) (0.2100,0.2500) (0.2107,0.2500) (0.2113,0.2500)
(0.2120,0.2500) (0.2127,0.2500) (0.2133,0.2500) (0.2140,0.2500) (0.2147,0.2500) (0.2153,0.2500)
(0.2160,0.2500) (0.2167,0.2500) (0.2173,0.2500) (0.2180,0.2500) (0.2187,0.2500) (0.2193,0.2500)
(0.2200,0.2500) (0.2207,0.2500) (0.2213,0.2500) (0.2220,0.2500) (0.2227,0.2571) (0.2233,0.2617)
(0.2240,0.2656) (0.2247,0.2656) (0.2253,0.2715) (0.2260,0.2764) (0.2267,0.2812) (0.2273,0.2812)
(0.2280,0.2812) (0.2287,0.2812) (0.2293,0.2812) (0.2300,0.2812) (0.2307,0.2852) (0.2313,0.2891)
(0.2320,0.2969) (0.2327,0.2969) (0.2333,0.3000) (0.2340,0.3047) (0.2347,0.3125) (0.2353,0.3125)
(0.2360,0.3125) (0.2367,0.3125) (0.2373,0.3125) (0.2380,0.3125) (0.2387,0.3125) (0.2393,0.3125)
(0.2400,0.3125) (0.2407,0.3125) (0.2413,0.3125) (0.2420,0.3125) (0.2427,0.3125) (0.2433,0.3125)
(0.2440,0.3125) (0.2447,0.3125) (0.2453,0.3125) (0.2460,0.3125) (0.2467,0.3125) (0.2473,0.3191)
(0.2480,0.3242) (0.2487,0.3281) (0.2493,0.3281) (0.2500,0.3333) (0.2507,0.3379) (0.2513,0.3438)
(0.2520,0.3438) (0.2527,0.3438) (0.2533,0.3438) (0.2540,0.3438) (0.2547,0.3438) (0.2553,0.3477)
(0.2560,0.3516) (0.2567,0.3594) (0.2573,0.3594) (0.2580,0.3618) (0.2587,0.3672) (0.2593,0.3750)
(0.2600,0.3750) (0.2607,0.3750) (0.2613,0.3750) (0.2620,0.3750) (0.2627,0.3750) (0.2633,0.3750)
(0.2640,0.3750) (0.2647,0.3750) (0.2653,0.3750) (0.2660,0.3750) (0.2667,0.3750) (0.2673,0.3750)
(0.2680,0.3750) (0.2687,0.3750) (0.2693,0.3750) (0.2700,0.3750) (0.2707,0.3750) (0.2713,0.3750)
(0.2720,0.3750) (0.2727,0.3750) (0.2733,0.3750) (0.2740,0.3750) (0.2747,0.3750) (0.2753,0.3750)
(0.2760,0.3750) (0.2767,0.3750) (0.2773,0.3750) (0.2780,0.3750) (0.2787,0.3750) (0.2793,0.3750)
(0.2800,0.3750) (0.2807,0.3750) (0.2813,0.3750) (0.2820,0.3750) (0.2827,0.3750) (0.2833,0.3750)
(0.2840,0.3750) (0.2847,0.3750) (0.2853,0.3750) (0.2860,0.3750) (0.2867,0.3750) (0.2873,0.3750)
(0.2880,0.3750) (0.2887,0.3750) (0.2893,0.3750) (0.2900,0.3750) (0.2907,0.3750) (0.2913,0.3750)
(0.2920,0.3750) (0.2927,0.3750) (0.2933,0.3750) (0.2940,0.3750) (0.2947,0.3750) (0.2953,0.3750)
(0.2960,0.3750) (0.2967,0.3809) (0.2973,0.3857) (0.2980,0.3906) (0.2987,0.3906) (0.2993,0.3945)
(0.3000,0.4000) (0.3007,0.4062) (0.3013,0.4062) (0.3020,0.4062) (0.3027,0.4062) (0.3033,0.4062)
(0.3040,0.4062) (0.3047,0.4095) (0.3053,0.4141) (0.3060,0.4219) (0.3067,0.4219) (0.3073,0.4238)
(0.3080,0.4297) (0.3087,0.4375) (0.3093,0.4375) (0.3100,0.4375) (0.3107,0.4375) (0.3113,0.4375)
(0.3120,0.4375) (0.3127,0.4375) (0.3133,0.4375) (0.3140,0.4375) (0.3147,0.4375) (0.3153,0.4375)
(0.3160,0.4375) (0.3167,0.4375) (0.3173,0.4375) (0.3180,0.4375) (0.3187,0.4375) (0.3193,0.4375)
(0.3200,0.4375) (0.3207,0.4375) (0.3213,0.4429) (0.3220,0.4473) (0.3227,0.4531) (0.3233,0.4531)
(0.3240,0.4570) (0.3247,0.4619) (0.3253,0.4688) (0.3260,0.4688) (0.3267,0.4688) (0.3273,0.4688)
(0.3280,0.4688) (0.3287,0.4688) (0.3293,0.4714) (0.3300,0.4766) (0.3307,0.4844) (0.3313,0.4844)
(0.3320,0.4857) (0.3327,0.4922) (0.3333,0.5000) (0.3340,0.5000) (0.3347,0.5000) (0.3353,0.5000)
(0.3360,0.5000) (0.3367,0.5000) (0.3373,0.5000) (0.3380,0.5000) (0.3387,0.5000) (0.3393,0.5000)
(0.3400,0.5000) (0.3407,0.5000) (0.3413,0.5000) (0.3420,0.5000) (0.3427,0.5000) (0.3433,0.5000)
(0.3440,0.5000) (0.3447,0.5000) (0.3453,0.5000) (0.3460,0.5000) (0.3467,0.5000) (0.3473,0.5000)
(0.3480,0.5000) (0.3487,0.5000) (0.3493,0.5000) (0.3500,0.5000) (0.3507,0.5000) (0.3513,0.5000)
(0.3520,0.5000) (0.3527,0.5000) (0.3533,0.5000) (0.3540,0.5000) (0.3547,0.5000) (0.3553,0.5000)
(0.3560,0.5000) (0.3567,0.5000) (0.3573,0.5000) (0.3580,0.5000) (0.3587,0.5000) (0.3593,0.5000)
(0.3600,0.5000) (0.3607,0.5000) (0.3613,0.5000) (0.3620,0.5000) (0.3627,0.5000) (0.3633,0.5000)
(0.3640,0.5000) (0.3647,0.5000) (0.3653,0.5000) (0.3660,0.5000) (0.3667,0.5000) (0.3673,0.5000)
(0.3680,0.5000) (0.3687,0.5000) (0.3693,0.5000) (0.3700,0.5000) (0.3707,0.5000) (0.3713,0.5000)
(0.3720,0.5000) (0.3727,0.5000) (0.3733,0.5000) (0.3740,0.5000) (0.3747,0.5000) (0.3753,0.5000)
(0.3760,0.5000) (0.3767,0.5000) (0.3773,0.5000) (0.3780,0.5000) (0.3787,0.5000) (0.3793,0.5000)
(0.3800,0.5000) (0.3807,0.5000) (0.3813,0.5000) (0.3820,0.5000) (0.3827,0.5000) (0.3833,0.5000)
(0.3840,0.5000) (0.3847,0.5000) (0.3853,0.5000) (0.3860,0.5000) (0.3867,0.5000) (0.3873,0.5000)
(0.3880,0.5000) (0.3887,0.5000) (0.3893,0.5000) (0.3900,0.5000) (0.3907,0.5000) (0.3913,0.5000)
(0.3920,0.5000) (0.3927,0.5000) (0.3933,0.5000) (0.3940,0.5000) (0.3947,0.5000) (0.3953,0.5000)
(0.3960,0.5000) (0.3967,0.5000) (0.3973,0.5000) (0.3980,0.5000) (0.3987,0.5000) (0.3993,0.5000)
(0.4000,0.5000) (0.4007,0.5000) (0.4013,0.5000) (0.4020,0.5000) (0.4027,0.5000) (0.4033,0.5000)
(0.4040,0.5000) (0.4047,0.5000) (0.4053,0.5000) (0.4060,0.5000) (0.4067,0.5000) (0.4073,0.5000)
(0.4080,0.5000) (0.4087,0.5000) (0.4093,0.5000) (0.4100,0.5000) (0.4107,0.5000) (0.4113,0.5000)
(0.4120,0.5000) (0.4127,0.5000) (0.4133,0.5000) (0.4140,0.5000) (0.4147,0.5000) (0.4153,0.5000)
(0.4160,0.5000) (0.4167,0.5000) (0.4173,0.5000) (0.4180,0.5000) (0.4187,0.5000) (0.4193,0.5000)
(0.4200,0.5000) (0.4207,0.5000) (0.4213,0.5000) (0.4220,0.5000) (0.4227,0.5000) (0.4233,0.5000)
(0.4240,0.5000) (0.4247,0.5000) (0.4253,0.5000) (0.4260,0.5000) (0.4267,0.5000) (0.4273,0.5000)
(0.4280,0.5000) (0.4287,0.5000) (0.4293,0.5000) (0.4300,0.5000) (0.4307,0.5000) (0.4313,0.5000)
(0.4320,0.5000) (0.4327,0.5000) (0.4333,0.5000) (0.4340,0.5000) (0.4347,0.5000) (0.4353,0.5000)
(0.4360,0.5000) (0.4367,0.5000) (0.4373,0.5000) (0.4380,0.5000) (0.4387,0.5000) (0.4393,0.5000)
(0.4400,0.5000) (0.4407,0.5000) (0.4413,0.5000) (0.4420,0.5000) (0.4427,0.5000) (0.4433,0.5000)
(0.4440,0.5000) (0.4447,0.5000) (0.4453,0.5000) (0.4460,0.5000) (0.4467,0.5000) (0.4473,0.5000)
(0.4480,0.5000) (0.4487,0.5000) (0.4493,0.5000) (0.4500,0.5000) (0.4507,0.5000) (0.4513,0.5000)
(0.4520,0.5000) (0.4527,0.5000) (0.4533,0.5000) (0.4540,0.5000) (0.4547,0.5000) (0.4553,0.5000)
(0.4560,0.5000) (0.4567,0.5000) (0.4573,0.5000) (0.4580,0.5000) (0.4587,0.5000) (0.4593,0.5000)
(0.4600,0.5000) (0.4607,0.5000) (0.4613,0.5000) (0.4620,0.5000) (0.4627,0.5000) (0.4633,0.5000)
(0.4640,0.5000) (0.4647,0.5000) (0.4653,0.5000) (0.4660,0.5000) (0.4667,0.5000) (0.4673,0.5000)
(0.4680,0.5000) (0.4687,0.5000) (0.4693,0.5000) (0.4700,0.5000) (0.4707,0.5000) (0.4713,0.5000)
(0.4720,0.5000) (0.4727,0.5000) (0.4733,0.5000) (0.4740,0.5000) (0.4747,0.5000) (0.4753,0.5000)
(0.4760,0.5000) (0.4767,0.5000) (0.4773,0.5000) (0.4780,0.5000) (0.4787,0.5000) (0.4793,0.5000)
(0.4800,0.5000) (0.4807,0.5000) (0.4813,0.5000) (0.4820,0.5000) (0.4827,0.5000) (0.4833,0.5000)
(0.4840,0.5000) (0.4847,0.5000) (0.4853,0.5000) (0.4860,0.5000) (0.4867,0.5000) (0.4873,0.5000)
(0.4880,0.5000) (0.4887,0.5000) (0.4893,0.5000) (0.4900,0.5000) (0.4907,0.5000) (0.4913,0.5000)
(0.4920,0.5000) (0.4927,0.5000) (0.4933,0.5000) (0.4940,0.5000) (0.4947,0.5000) (0.4953,0.5000)
(0.4960,0.5000) (0.4967,0.5000) (0.4973,0.5000) (0.4980,0.5000) (0.4987,0.5000) (0.4993,0.5000)
(0.5000,0.5000) (0.5007,0.5000) (0.5013,0.5000) (0.5020,0.5000) (0.5027,0.5000) (0.5033,0.5000)
(0.5040,0.5000) (0.5047,0.5000) (0.5053,0.5000) (0.5060,0.5000) (0.5067,0.5000) (0.5073,0.5000)
(0.5080,0.5000) (0.5087,0.5000) (0.5093,0.5000) (0.5100,0.5000) (0.5107,0.5000) (0.5113,0.5000)
(0.5120,0.5000) (0.5127,0.5000) (0.5133,0.5000) (0.5140,0.5000) (0.5147,0.5000) (0.5153,0.5000)
(0.5160,0.5000) (0.5167,0.5000) (0.5173,0.5000) (0.5180,0.5000) (0.5187,0.5000) (0.5193,0.5000)
(0.5200,0.5000) (0.5207,0.5000) (0.5213,0.5000) (0.5220,0.5000) (0.5227,0.5000) (0.5233,0.5000)
(0.5240,0.5000) (0.5247,0.5000) (0.5253,0.5000) (0.5260,0.5000) (0.5267,0.5000) (0.5273,0.5000)
(0.5280,0.5000) (0.5287,0.5000) (0.5293,0.5000) (0.5300,0.5000) (0.5307,0.5000) (0.5313,0.5000)
(0.5320,0.5000) (0.5327,0.5000) (0.5333,0.5000) (0.5340,0.5000) (0.5347,0.5000) (0.5353,0.5000)
(0.5360,0.5000) (0.5367,0.5000) (0.5373,0.5000) (0.5380,0.5000) (0.5387,0.5000) (0.5393,0.5000)
(0.5400,0.5000) (0.5407,0.5000) (0.5413,0.5000) (0.5420,0.5000) (0.5427,0.5000) (0.5433,0.5000)
(0.5440,0.5000) (0.5447,0.5000) (0.5453,0.5000) (0.5460,0.5000) (0.5467,0.5000) (0.5473,0.5000)
(0.5480,0.5000) (0.5487,0.5000) (0.5493,0.5000) (0.5500,0.5000) (0.5507,0.5000) (0.5513,0.5000)
(0.5520,0.5000) (0.5527,0.5000) (0.5533,0.5000) (0.5540,0.5000) (0.5547,0.5000) (0.5553,0.5000)
(0.5560,0.5000) (0.5567,0.5000) (0.5573,0.5000) (0.5580,0.5000) (0.5587,0.5000) (0.5593,0.5000)
(0.5600,0.5000) (0.5607,0.5000) (0.5613,0.5000) (0.5620,0.5000) (0.5627,0.5000) (0.5633,0.5000)
(0.5640,0.5000) (0.5647,0.5000) (0.5653,0.5000) (0.5660,0.5000) (0.5667,0.5000) (0.5673,0.5000)
(0.5680,0.5000) (0.5687,0.5000) (0.5693,0.5000) (0.5700,0.5000) (0.5707,0.5000) (0.5713,0.5000)
(0.5720,0.5000) (0.5727,0.5000) (0.5733,0.5000) (0.5740,0.5000) (0.5747,0.5000) (0.5753,0.5000)
(0.5760,0.5000) (0.5767,0.5000) (0.5773,0.5000) (0.5780,0.5000) (0.5787,0.5000) (0.5793,0.5000)
(0.5800,0.5000) (0.5807,0.5000) (0.5813,0.5000) (0.5820,0.5000) (0.5827,0.5000) (0.5833,0.5000)
(0.5840,0.5000) (0.5847,0.5000) (0.5853,0.5000) (0.5860,0.5000) (0.5867,0.5000) (0.5873,0.5000)
(0.5880,0.5000) (0.5887,0.5000) (0.5893,0.5000) (0.5900,0.5000) (0.5907,0.5000) (0.5913,0.5000)
(0.5920,0.5000) (0.5927,0.5000) (0.5933,0.5000) (0.5940,0.5000) (0.5947,0.5000) (0.5953,0.5000)
(0.5960,0.5000) (0.5967,0.5000) (0.5973,0.5000) (0.5980,0.5000) (0.5987,0.5000) (0.5993,0.5000)
(0.6000,0.5000) (0.6007,0.5000) (0.6013,0.5000) (0.6020,0.5000) (0.6027,0.5000) (0.6033,0.5000)
(0.6040,0.5000) (0.6047,0.5000) (0.6053,0.5000) (0.6060,0.5000) (0.6067,0.5000) (0.6073,0.5000)
(0.6080,0.5000) (0.6087,0.5000) (0.6093,0.5000) (0.6100,0.5000) (0.6107,0.5000) (0.6113,0.5000)
(0.6120,0.5000) (0.6127,0.5000) (0.6133,0.5000) (0.6140,0.5000) (0.6147,0.5000) (0.6153,0.5000)
(0.6160,0.5000) (0.6167,0.5000) (0.6173,0.5000) (0.6180,0.5000) (0.6187,0.5000) (0.6193,0.5000)
(0.6200,0.5000) (0.6207,0.5000) (0.6213,0.5000) (0.6220,0.5000) (0.6227,0.5000) (0.6233,0.5000)
(0.6240,0.5000) (0.6247,0.5000) (0.6253,0.5000) (0.6260,0.5000) (0.6267,0.5000) (0.6273,0.5000)
(0.6280,0.5000) (0.6287,0.5000) (0.6293,0.5000) (0.6300,0.5000) (0.6307,0.5000) (0.6313,0.5000)
(0.6320,0.5000) (0.6327,0.5000) (0.6333,0.5000) (0.6340,0.5000) (0.6347,0.5000) (0.6353,0.5000)
(0.6360,0.5000) (0.6367,0.5000) (0.6373,0.5000) (0.6380,0.5000) (0.6387,0.5000) (0.6393,0.5000)
(0.6400,0.5000) (0.6407,0.5000) (0.6413,0.5000) (0.6420,0.5000) (0.6427,0.5000) (0.6433,0.5000)
(0.6440,0.5000) (0.6447,0.5000) (0.6453,0.5000) (0.6460,0.5000) (0.6467,0.5000) (0.6473,0.5000)
(0.6480,0.5000) (0.6487,0.5000) (0.6493,0.5000) (0.6500,0.5000) (0.6507,0.5000) (0.6513,0.5000)
(0.6520,0.5000) (0.6527,0.5000) (0.6533,0.5000) (0.6540,0.5000) (0.6547,0.5000) (0.6553,0.5000)
(0.6560,0.5000) (0.6567,0.5000) (0.6573,0.5000) (0.6580,0.5000) (0.6587,0.5000) (0.6593,0.5000)
(0.6600,0.5000) (0.6607,0.5000) (0.6613,0.5000) (0.6620,0.5000) (0.6627,0.5000) (0.6633,0.5000)
(0.6640,0.5000) (0.6647,0.5000) (0.6653,0.5000) (0.6660,0.5000) (0.6667,0.5000) (0.6673,0.5078)
(0.6680,0.5143) (0.6687,0.5156) (0.6693,0.5156) (0.6700,0.5234) (0.6707,0.5286) (0.6713,0.5312)
(0.6720,0.5312) (0.6727,0.5312) (0.6733,0.5312) (0.6740,0.5312) (0.6747,0.5312) (0.6753,0.5381)
(0.6760,0.5430) (0.6767,0.5469) (0.6773,0.5469) (0.6780,0.5527) (0.6787,0.5571) (0.6793,0.5625)
(0.6800,0.5625) (0.6807,0.5625) (0.6813,0.5625) (0.6820,0.5625) (0.6827,0.5625) (0.6833,0.5625)
(0.6840,0.5625) (0.6847,0.5625) (0.6853,0.5625) (0.6860,0.5625) (0.6867,0.5625) (0.6873,0.5625)
(0.6880,0.5625) (0.6887,0.5625) (0.6893,0.5625) (0.6900,0.5625) (0.6907,0.5625) (0.6913,0.5625)
(0.6920,0.5703) (0.6927,0.5762) (0.6933,0.5781) (0.6940,0.5781) (0.6947,0.5859) (0.6953,0.5905)
(0.6960,0.5938) (0.6967,0.5938) (0.6973,0.5938) (0.6980,0.5938) (0.6987,0.5938) (0.6993,0.5938)
(0.7000,0.6000) (0.7007,0.6055) (0.7013,0.6094) (0.7020,0.6094) (0.7027,0.6143) (0.7033,0.6191)
(0.7040,0.6250) (0.7047,0.6250) (0.7053,0.6250) (0.7060,0.6250) (0.7067,0.6250) (0.7073,0.6250)
(0.7080,0.6250) (0.7087,0.6250) (0.7093,0.6250) (0.7100,0.6250) (0.7107,0.6250) (0.7113,0.6250)
(0.7120,0.6250) (0.7127,0.6250) (0.7133,0.6250) (0.7140,0.6250) (0.7147,0.6250) (0.7153,0.6250)
(0.7160,0.6250) (0.7167,0.6250) (0.7173,0.6250) (0.7180,0.6250) (0.7187,0.6250) (0.7193,0.6250)
(0.7200,0.6250) (0.7207,0.6250) (0.7213,0.6250) (0.7220,0.6250) (0.7227,0.6250) (0.7233,0.6250)
(0.7240,0.6250) (0.7247,0.6250) (0.7253,0.6250) (0.7260,0.6250) (0.7267,0.6250) (0.7273,0.6250)
(0.7280,0.6250) (0.7287,0.6250) (0.7293,0.6250) (0.7300,0.6250) (0.7307,0.6250) (0.7313,0.6250)
(0.7320,0.6250) (0.7327,0.6250) (0.7333,0.6250) (0.7340,0.6250) (0.7347,0.6250) (0.7353,0.6250)
(0.7360,0.6250) (0.7367,0.6250) (0.7373,0.6250) (0.7380,0.6250) (0.7387,0.6250) (0.7393,0.6250)
(0.7400,0.6250) (0.7407,0.6250) (0.7413,0.6328) (0.7420,0.6382) (0.7427,0.6406) (0.7433,0.6406)
(0.7440,0.6484) (0.7447,0.6523) (0.7453,0.6562) (0.7460,0.6562) (0.7467,0.6562) (0.7473,0.6562)
(0.7480,0.6562) (0.7487,0.6562) (0.7493,0.6621) (0.7500,0.6667) (0.7507,0.6719) (0.7513,0.6719)
(0.7520,0.6758) (0.7527,0.6809) (0.7533,0.6875) (0.7540,0.6875) (0.7547,0.6875) (0.7553,0.6875)
(0.7560,0.6875) (0.7567,0.6875) (0.7573,0.6875) (0.7580,0.6875) (0.7587,0.6875) (0.7593,0.6875)
(0.7600,0.6875) (0.7607,0.6875) (0.7613,0.6875) (0.7620,0.6875) (0.7627,0.6875) (0.7633,0.6875)
(0.7640,0.6875) (0.7647,0.6875) (0.7653,0.6875) (0.7660,0.6953) (0.7667,0.7000) (0.7673,0.7031)
(0.7680,0.7031) (0.7687,0.7109) (0.7693,0.7148) (0.7700,0.7188) (0.7707,0.7188) (0.7713,0.7188)
(0.7720,0.7188) (0.7727,0.7188) (0.7733,0.7188) (0.7740,0.7236) (0.7747,0.7285) (0.7753,0.7344)
(0.7760,0.7344) (0.7767,0.7383) (0.7773,0.7429) (0.7780,0.7500) (0.7787,0.7500) (0.7793,0.7500)
(0.7800,0.7500) (0.7807,0.7500) (0.7813,0.7500) (0.7820,0.7500) (0.7827,0.7500) (0.7833,0.7500)
(0.7840,0.7500) (0.7847,0.7500) (0.7853,0.7500) (0.7860,0.7500) (0.7867,0.7500) (0.7873,0.7500)
(0.7880,0.7500) (0.7887,0.7500) (0.7893,0.7500) (0.7900,0.7500) (0.7907,0.7500) (0.7913,0.7500)
(0.7920,0.7500) (0.7927,0.7500) (0.7933,0.7500) (0.7940,0.7500) (0.7947,0.7500) (0.7953,0.7500)
(0.7960,0.7500) (0.7967,0.7500) (0.7973,0.7500) (0.7980,0.7500) (0.7987,0.7500) (0.7993,0.7500)
(0.8000,0.7500) (0.8007,0.7500) (0.8013,0.7500) (0.8020,0.7500) (0.8027,0.7500) (0.8033,0.7500)
(0.8040,0.7500) (0.8047,0.7500) (0.8053,0.7500) (0.8060,0.7500) (0.8067,0.7500) (0.8073,0.7500)
(0.8080,0.7500) (0.8087,0.7500) (0.8093,0.7500) (0.8100,0.7500) (0.8107,0.7500) (0.8113,0.7500)
(0.8120,0.7500) (0.8127,0.7500) (0.8133,0.7500) (0.8140,0.7500) (0.8147,0.7500) (0.8153,0.7500)
(0.8160,0.7500) (0.8167,0.7500) (0.8173,0.7500) (0.8180,0.7500) (0.8187,0.7500) (0.8193,0.7500)
(0.8200,0.7500) (0.8207,0.7500) (0.8213,0.7500) (0.8220,0.7500) (0.8227,0.7500) (0.8233,0.7500)
(0.8240,0.7500) (0.8247,0.7500) (0.8253,0.7500) (0.8260,0.7500) (0.8267,0.7500) (0.8273,0.7500)
(0.8280,0.7500) (0.8287,0.7500) (0.8293,0.7500) (0.8300,0.7500) (0.8307,0.7500) (0.8313,0.7500)
(0.8320,0.7500) (0.8327,0.7500) (0.8333,0.7500) (0.8340,0.7500) (0.8347,0.7500) (0.8353,0.7500)
(0.8360,0.7500) (0.8367,0.7500) (0.8373,0.7500) (0.8380,0.7500) (0.8387,0.7500) (0.8393,0.7500)
(0.8400,0.7500) (0.8407,0.7500) (0.8413,0.7500) (0.8420,0.7500) (0.8427,0.7500) (0.8433,0.7500)
(0.8440,0.7500) (0.8447,0.7500) (0.8453,0.7500) (0.8460,0.7500) (0.8467,0.7500) (0.8473,0.7500)
(0.8480,0.7500) (0.8487,0.7500) (0.8493,0.7500) (0.8500,0.7500) (0.8507,0.7500) (0.8513,0.7500)
(0.8520,0.7500) (0.8527,0.7500) (0.8533,0.7500) (0.8540,0.7500) (0.8547,0.7500) (0.8553,0.7500)
(0.8560,0.7500) (0.8567,0.7500) (0.8573,0.7500) (0.8580,0.7500) (0.8587,0.7500) (0.8593,0.7500)
(0.8600,0.7500) (0.8607,0.7500) (0.8613,0.7500) (0.8620,0.7500) (0.8627,0.7500) (0.8633,0.7500)
(0.8640,0.7500) (0.8647,0.7500) (0.8653,0.7500) (0.8660,0.7500) (0.8667,0.7500) (0.8673,0.7500)
(0.8680,0.7500) (0.8687,0.7500) (0.8693,0.7500) (0.8700,0.7500) (0.8707,0.7500) (0.8713,0.7500)
(0.8720,0.7500) (0.8727,0.7500) (0.8733,0.7500) (0.8740,0.7500) (0.8747,0.7500) (0.8753,0.7500)
(0.8760,0.7500) (0.8767,0.7500) (0.8773,0.7500) (0.8780,0.7500) (0.8787,0.7500) (0.8793,0.7500)
(0.8800,0.7500) (0.8807,0.7500) (0.8813,0.7500) (0.8820,0.7500) (0.8827,0.7500) (0.8833,0.7500)
(0.8840,0.7500) (0.8847,0.7500) (0.8853,0.7500) (0.8860,0.7500) (0.8867,0.7500) (0.8873,0.7500)
(0.8880,0.7500) (0.8887,0.7500) (0.8893,0.7571) (0.8900,0.7617) (0.8907,0.7656) (0.8913,0.7656)
(0.8920,0.7715) (0.8927,0.7764) (0.8933,0.7812) (0.8940,0.7812) (0.8947,0.7812) (0.8953,0.7812)
(0.8960,0.7812) (0.8967,0.7812) (0.8973,0.7852) (0.8980,0.7891) (0.8987,0.7969) (0.8993,0.7969)
(0.9000,0.8000) (0.9007,0.8047) (0.9013,0.8125) (0.9020,0.8125) (0.9027,0.8125) (0.9033,0.8125)
(0.9040,0.8125) (0.9047,0.8125) (0.9053,0.8125) (0.9060,0.8125) (0.9067,0.8125) (0.9073,0.8125)
(0.9080,0.8125) (0.9087,0.8125) (0.9093,0.8125) (0.9100,0.8125) (0.9107,0.8125) (0.9113,0.8125)
(0.9120,0.8125) (0.9127,0.8125) (0.9133,0.8125) (0.9140,0.8191) (0.9147,0.8242) (0.9153,0.8281)
(0.9160,0.8281) (0.9167,0.8333) (0.9173,0.8379) (0.9180,0.8438) (0.9187,0.8438) (0.9193,0.8438)
(0.9200,0.8438) (0.9207,0.8438) (0.9213,0.8438) (0.9220,0.8477) (0.9227,0.8516) (0.9233,0.8594)
(0.9240,0.8594) (0.9247,0.8618) (0.9253,0.8672) (0.9260,0.8750) (0.9267,0.8750) (0.9273,0.8750)
(0.9280,0.8750) (0.9287,0.8750) (0.9293,0.8750) (0.9300,0.8750) (0.9307,0.8750) (0.9313,0.8750)
(0.9320,0.8750) (0.9327,0.8750) (0.9333,0.8750) (0.9340,0.8750) (0.9347,0.8750) (0.9353,0.8750)
(0.9360,0.8750) (0.9367,0.8750) (0.9373,0.8750) (0.9380,0.8750) (0.9387,0.8750) (0.9393,0.8750)
(0.9400,0.8750) (0.9407,0.8750) (0.9413,0.8750) (0.9420,0.8750) (0.9427,0.8750) (0.9433,0.8750)
(0.9440,0.8750) (0.9447,0.8750) (0.9453,0.8750) (0.9460,0.8750) (0.9467,0.8750) (0.9473,0.8750)
(0.9480,0.8750) (0.9487,0.8750) (0.9493,0.8750) (0.9500,0.8750) (0.9507,0.8750) (0.9513,0.8750)
(0.9520,0.8750) (0.9527,0.8750) (0.9533,0.8750) (0.9540,0.8750) (0.9547,0.8750) (0.9553,0.8750)
(0.9560,0.8750) (0.9567,0.8750) (0.9573,0.8750) (0.9580,0.8750) (0.9587,0.8750) (0.9593,0.8750)
(0.9600,0.8750) (0.9607,0.8750) (0.9613,0.8750) (0.9620,0.8750) (0.9627,0.8750) (0.9633,0.8809)
(0.9640,0.8857) (0.9647,0.8906) (0.9653,0.8906) (0.9660,0.8945) (0.9667,0.9000) (0.9673,0.9062)
(0.9680,0.9062) (0.9687,0.9062) (0.9693,0.9062) (0.9700,0.9062) (0.9707,0.9062) (0.9713,0.9095)
(0.9720,0.9141) (0.9727,0.9219) (0.9733,0.9219) (0.9740,0.9238) (0.9747,0.9297) (0.9753,0.9375)
(0.9760,0.9375) (0.9767,0.9375) (0.9773,0.9375) (0.9780,0.9375) (0.9787,0.9375) (0.9793,0.9375)
(0.9800,0.9375) (0.9807,0.9375) (0.9813,0.9375) (0.9820,0.9375) (0.9827,0.9375) (0.9833,0.9375)
(0.9840,0.9375) (0.9847,0.9375) (0.9853,0.9375) (0.9860,0.9375) (0.9867,0.9375) (0.9873,0.9375)
(0.9880,0.9429) (0.9887,0.9473) (0.9893,0.9531) (0.9900,0.9531) (0.9907,0.9570) (0.9913,0.9619)
(0.9920,0.9688) (0.9927,0.9688) (0.9933,0.9688) (0.9940,0.9688) (0.9947,0.9688) (0.9953,0.9688)
(0.9960,0.9714) (0.9967,0.9766) (0.9973,0.9844) (0.9980,0.9844) (0.9987,0.9857) (0.9993,0.9922)
(1.0000,1.0000)
};

\end{tikzpicture}
\caption{L'escalier de Cantor : continu, croissant, constant sur chaque « marche » (intervalle
contigu de $K$), donc $c'=0$ \pp, et pourtant $c(0)=0$, $c(1)=1$.}
\label{fig:cantor}
\end{figure}

\begin{proposition}\label{prop:cantor}
La fonction $c$ vérifie :
\begin{enumerate}[label=\textup{(\alph*)},leftmargin=2.2em,itemsep=1pt]
  \item $c$ est croissante, $c(0)=0$, $c(1)=1$, et $c$ est \emph{surjective} sur $[0,1]$ ;
  \item $c$ est \emph{continue} ;
  \item $c'(x)=0$ pour tout $x\in[0,1]\setminus K$, donc pour presque tout $x$, et
  $\displaystyle\int_0^1 c'(x)\dx=0\neq1=c(1)-c(0)$.
\end{enumerate}
\end{proposition}

\begin{proof}
(a) \emph{Croissance.} Soient $x<y$ dans $[0,1]$. Choisissons leurs développements triadiques
« propres ». Si $x<y$, au premier rang $i_0$ où leurs chiffres diffèrent, on a $a_{i_0}(x)<a_{i_0}(y)$.
L'application chiffre à chiffre $a\mapsto b$ de la définition~\ref{def:cantor} est croissante et,
au rang $i_0$, donne $b_{i_0}(x)\le b_{i_0}(y)$, l'égalité n'ayant lieu que si un $a=1$ intervient ;
un examen des cas (en tenant compte de la troncature au premier chiffre $1$) montre dans tous les cas
$c(x)\le c(y)$.
\emph{Surjectivité.} Soit $t\in[0,1]$, de développement binaire $t=\sum_i \beta_i2^{-i}$,
$\beta_i\in\{0,1\}$. Posons $x=\sum_i(2\beta_i)3^{-i}\in K$ ; alors, par construction, $c(x)=t$. Enfin
$c(0)=0$ et $c(1)=c(0{.}222\ldots_3)=0{.}111\ldots_2=1$.

(b) \emph{Continuité.} $c$ est croissante et son image contient $[0,1]$ tout entier (surjectivité).
Une fonction croissante ne peut présenter qu'au plus des sauts, et un saut en $x_0$ priverait l'image
de l'intervalle ouvert $\bigl(c(x_0^-),c(x_0^+)\bigr)$ ; l'image ne serait pas $[0,1]$. Donc $c$ est
sans saut, c'est-à-dire continue.

(c) Chaque intervalle contigu de $K$ est un intervalle ouvert sur lequel $c$ est constante, donc
$c'=0$ y est identiquement nulle. Comme $\lambda([0,1]\setminus K)=1$, on a $c'=0$ \pp, d'où
$\int_0^1 c'=0$. Or $c(1)-c(0)=1$.
\end{proof}

\begin{remarque}
$c$ est höldérienne d'exposant $\alpha=\log2/\log3\approx0{,}63$ : sur deux points d'un même intervalle
triadique de niveau $n$ (longueur $3^{-n}$), $c$ varie d'au plus $2^{-n}=(3^{-n})^{\alpha}$, d'où
$\abs{c(x)-c(y)}\le C\abs{x-y}^{\alpha}$. Elle n'est lipschitzienne sur aucun voisinage d'un point de
$K$ (sinon on aurait $\alpha\ge1$).
\end{remarque}

\begin{remarque}[lecture par les mesures]
À $c$ croissante on associe la mesure de Lebesgue--Stieltjes $\mu_c$ définie par
$\mu_c(]x,y])=c(y)-c(x)$. C'est une probabilité, \emph{sans atome} (car $c$ continue), \emph{portée
par $K$} (car $c$ constante hors de $K$), donc \emph{étrangère} à $\lambda$ : $\mu_c\perp\lambda$.
L'égalité $\int_a^b F'=F(b)-F(a)$ équivaut à $d\mu_F=F'\,d\lambda$, c'est-à-dire à l'absolue
continuité de $\mu_F$ par rapport à $\lambda$ ; l'escalier échoue précisément parce que $\mu_c$ est
singulière. C'est le diagnostic qui guide toute la suite.
\end{remarque}

% ============================================================
\section{Fonctions à variation bornée}
% ============================================================

\begin{definition}
La \emph{variation totale} de $f$ sur $[a,b]$ est
$\Var_a^b(f)=\sup\sum_{k}\abs{f(x_k)-f(x_{k-1})}$, le sup portant sur toutes les subdivisions
$a=x_0<\cdots<x_n=b$. On note $f\in BV[a,b]$ si $\Var_a^b(f)<\infty$.
\end{definition}

On vérifie aisément l'\emph{additivité} $\Var_a^c(f)=\Var_a^x(f)+\Var_x^c(f)$ pour $a\le x\le c$
(toute subdivision se raffine en insérant $x$).

\begin{theoreme}[Jordan]\label{thm:jordan}
$f\in BV[a,b]$ si et seulement si $f=g-h$ avec $g,h$ croissantes.
\end{theoreme}

\begin{proof}
$(\Leftarrow)$ Une fonction croissante $\varphi$ est $BV$ avec $\Var_a^b(\varphi)=\varphi(b)-\varphi(a)$
(toute somme télescope). Et $BV$ est un espace vectoriel : $\Var(f_1+f_2)\le\Var(f_1)+\Var(f_2)$ par
inégalité triangulaire. Donc $g-h\in BV$.

$(\Rightarrow)$ Posons $v(x)=\Var_a^x(f)$, croissante par additivité. Montrons que $v-f$ est
croissante : pour $x<y$,
\[
(v-f)(y)-(v-f)(x)=\underbrace{\Var_x^y(f)}_{\ge\,\abs{f(y)-f(x)}}-\bigl(f(y)-f(x)\bigr)\ge
\abs{f(y)-f(x)}-\bigl(f(y)-f(x)\bigr)\ge0.
\]
Ainsi $f=v-(v-f)$ est différence de deux fonctions croissantes.
\end{proof}

\begin{corollaire}\label{cor:bv}
Toute $f\in BV[a,b]$ est dérivable \pp, $f'\in L^1$, et $\int_a^b\abs{f'}\le\Var_a^b(f)$.
\end{corollaire}

\begin{proof}
Écrire $f=g-h$ (théorème~\ref{thm:jordan}) et appliquer le théorème~\ref{thm:mono} à $g$ et $h$ :
$g',h'$ existent \pp, sont dans $L^1$, donc $f'=g'-h'$ aussi. Enfin $\abs{f'}\le g'+h'$ et
$\int(g'+h')\le\bigl(g(b)-g(a)\bigr)+\bigl(h(b)-h(a)\bigr)=\Var_a^b(f)$ (avec $g=v$, $h=v-f$).
\end{proof}

L'escalier de Cantor, croissant, est dans $BV$ (avec $\Var_0^1(c)=1$) et fournit déjà le
contre-exemple : \emph{la variation bornée ne suffit pas} au TFA. Il faut restreindre davantage.

% ============================================================
\section{Continuité absolue, et deux lemmes géométriques}
% ============================================================

\begin{definition}\label{def:ac}
$F:[a,b]\to\R$ est \emph{absolument continue} ($F\in AC[a,b]$) si : pour tout $\eps>0$ il existe
$\delta>0$ tel que, pour toute famille finie d'intervalles ouverts disjoints $(a_k,b_k)_{k=1}^n$,
\[
\sum_{k=1}^n(b_k-a_k)<\delta\ \Longrightarrow\ \sum_{k=1}^n\abs{F(b_k)-F(a_k)}<\eps.
\]
\end{definition}

\begin{remarque}
La même implication vaut alors pour une famille \emph{dénombrable} d'intervalles disjoints de longueur
totale $<\delta$ : appliquer la condition aux sommes partielles finies (chacune $<\eps$) et passer à
la limite. On l'utilisera librement.
\end{remarque}

\begin{proposition}[hiérarchie]\label{prop:hier}
Sur $[a,b]$ : $\Lip\subset AC\subset BV$, et $AC$ entraîne la continuité uniforme.
\end{proposition}

\begin{proof}
\emph{$\Lip\subset AC$ :} si $\abs{F(x)-F(y)}\le L\abs{x-y}$, alors
$\sum\abs{F(b_k)-F(a_k)}\le L\sum(b_k-a_k)$ ; prendre $\delta=\eps/L$.
\emph{$AC\Rightarrow$ continuité uniforme :} c'est le cas $n=1$ de la définition.
\emph{$AC\subset BV$ :} soit $\delta$ associé à $\eps=1$. Découpons $[a,b]$ en $M=\lceil(b-a)/\delta\rceil$
sous-intervalles $J_1,\dots,J_M$ de longueur $<\delta$. Sur chaque $J_m$, toute subdivision est une
famille d'intervalles disjoints de longueur totale $<\delta$, donc $\Var_{J_m}(F)\le1$. Par additivité,
$\Var_a^b(F)\le M<\infty$.
\end{proof}

\begin{exemple}\label{ex:integrale-ac}
Si $g\in L^1$, alors $F(x)=\int_a^x g$ est $AC$. En effet l'absolue continuité de l'intégrale de
Lebesgue donne, pour $\sum(b_k-a_k)<\delta$ bien choisi,
$\sum\abs{F(b_k)-F(a_k)}\le\int_{\bigcup(a_k,b_k)}\abs g<\eps$.
\end{exemple}

\begin{contreexemple}\label{ce:cantor-nonac}
L'escalier $c$ est continu et $BV$, mais \emph{pas} $AC$. À tout niveau $n$, les $2^n$ intervalles de
$K_n$ ont une longueur totale $(2/3)^n\to0$, alors que $c$ y accomplit \emph{toute} sa croissance : la
somme des accroissements de $c$ sur ces intervalles vaut $1$ (leurs images pavent $[0,1]$). Aucun
$\delta$ ne convient pour $\eps=\tfrac12$.
\end{contreexemple}

On établit maintenant deux lemmes géométriques, cœur technique de tout ce qui suit.

\begin{lemme}[image d'un ensemble à dérivée nulle]\label{lem:image}
Soit $H:[a,b]\to\R$ et $A\subset[a,b]$ tels que $H$ soit dérivable en chaque point de $A$ avec
$H'\equiv0$ sur $A$. Alors $\lstar\bigl(H(A)\bigr)=0$.
\end{lemme}

\begin{proof}
Fixons $\eps>0$ et, pour $n\ge1$, posons
\[
A_n=\Bigl\{x\in A:\ \abs{H(y)-H(x)}\le\eps\abs{y-x}\ \text{ pour tout }y\text{ avec }\abs{y-x}\le\tfrac1n\Bigr\}.
\]
Comme $H'(x)=0$ sur $A$, tout $x\in A$ appartient à $A_n$ pour $n$ assez grand : $A=\bigcup_n A_n$,
suite croissante. Fixons $n$ et découpons $[a,b]$ en cellules $J$ de longueur $\le1/n$. Si
$x,x'\in A_n\cap J$, alors $\abs{x-x'}\le\lambda(J)\le1/n$, donc $\abs{H(x)-H(x')}\le\eps\abs{x-x'}\le
\eps\lambda(J)$. Ainsi $\diam H(A_n\cap J)\le\eps\lambda(J)$, d'où $\lstar\bigl(H(A_n\cap J)\bigr)\le
\eps\lambda(J)$. En sommant sur les cellules recouvrant $[a,b]$ (de longueur totale $b-a$) :
\[
\lstar\bigl(H(A_n)\bigr)\le\sum_J\eps\lambda(J)=\eps(b-a).
\]
Comme $A_n\uparrow A$, on a $H(A_n)\uparrow H(A)$, et par continuité croissante de $\lstar$ (proposition~\ref{prop:contcroiss})
$\lstar\bigl(H(A)\bigr)=\lim_n\lstar\bigl(H(A_n)\bigr)\le\eps(b-a)$. Enfin $\eps\to0$ donne
$\lstar\bigl(H(A)\bigr)=0$.
\end{proof}

\begin{proposition}[$AC$ vérifie la propriété (N) de Luzin]\label{prop:luzin}
Si $F\in AC[a,b]$ et $\lambda(Z)=0$, alors $\lambda\bigl(F(Z)\bigr)=0$.
\end{proposition}

\begin{proof}
Fixons $\eps>0$ et $\delta>0$ associé par l'absolue continuité. Comme $\lambda(Z)=0$, la régularité
extérieure (proposition~\ref{prop:regext}) fournit un ouvert $U\supset Z$ avec $\lambda(U)<\delta$ ;
écrivons $U\cap(a,b)=\bigsqcup_k(a_k,b_k)$ (intervalles
disjoints), de sorte que $\sum_k(b_k-a_k)\le\lambda(U)<\delta$. Sur chaque $[a_k,b_k]$, $F$ continue
atteint son max en $p_k$ et son min en $q_k$ ; alors $F\bigl(Z\cap(a_k,b_k)\bigr)\subset[F(q_k),F(p_k)]$,
de longueur $\abs{F(p_k)-F(q_k)}$. Les couples $\{p_k,q_k\}$ définissent des intervalles deux à deux
disjoints (les $(a_k,b_k)$ le sont), de longueur totale $<\delta$. Par absolue continuité (version
dénombrable),
\[
\lstar\bigl(F(Z)\bigr)\le\sum_k\abs{F(p_k)-F(q_k)}\le\eps.
\]
Comme $\eps>0$ est arbitraire, $\lambda\bigl(F(Z)\bigr)=0$.
\end{proof}

\begin{lemme}[rigidité]\label{lem:rigidite}
Si $F\in AC[a,b]$ et $F'=0$ \pp, alors $F$ est constante.
\end{lemme}

\begin{proof}
Soit $A=\{x:F'(x)=0\}$, de mesure pleine : $\lambda([a,b]\setminus A)=0$. Le lemme~\ref{lem:image}
donne $\lambda\bigl(F(A)\bigr)=0$, et la proposition~\ref{prop:luzin} donne
$\lambda\bigl(F([a,b]\setminus A)\bigr)=0$. Donc $\lambda\bigl(F([a,b])\bigr)=0$. Mais $F$ est continue
sur l'intervalle $[a,b]$, donc $F([a,b])$ est un intervalle (théorème des valeurs intermédiaires) ;
un intervalle de mesure nulle est réduit à un point. Donc $F$ est constante.
\end{proof}

\begin{remarque}
Le lemme est \emph{faux} sans $AC$ : l'escalier de Cantor est continu, $F'=0$ \pp, mais non constant.
Ici, c'est la proposition~\ref{prop:luzin} qui tombe en défaut --- Cantor viole la propriété (N),
puisqu'il envoie $K$ (négligeable) sur $[0,1]$ (de mesure $1$). C'est le seul ingrédient qui lui
manque pour être $AC$.
\end{remarque}

% ============================================================
\section{Volet I : dérivation de l'intégrale}
% ============================================================

On démontre ici que dériver l'intégrale redonne l'intégrande. Un lemme d'annulation est d'abord requis.

\begin{lemme}[intégrale nulle]\label{lem:nul}
Si $g\in L^1$ et $\int_a^x g=0$ pour tout $x\in[a,b]$, alors $g=0$ \pp.
\end{lemme}

\begin{proof}
Par différence, $\int_c^d g=0$ pour tout sous-intervalle $[c,d]$. Tout ouvert $U\subset(a,b)$ étant
réunion dénombrable disjointe d'intervalles, la convergence dominée donne $\int_U g=0$. Par régularité extérieure (proposition~\ref{prop:regext})
extérieure, pour tout mesurable $E$ et tout $\eps>0$ il existe un ouvert $U\supset E$ avec
$\lambda(U\setminus E)<\eps$, d'où $\abs{\int_E g}=\abs{\int_U g-\int_{U\setminus E}g}\le
\int_{U\setminus E}\abs g\to0$ ; donc $\int_E g=0$ pour tout mesurable $E$. En prenant $E=\{g>0\}$ on
obtient $\int_{\{g>0\}}g=0$ avec $g>0$ sur $E$, donc $\lambda(\{g>0\})=0$ ; de même
$\lambda(\{g<0\})=0$. Ainsi $g=0$ \pp.
\end{proof}

\begin{theoreme}[Lebesgue --- dérivation de l'intégrale]\label{thm:volet1}
Soit $g\in L^1([a,b])$ et $G(x)=\int_a^x g$. Alors $G\in AC$, $G$ est dérivable \pp, et $G'=g$ \pp.
\end{theoreme}

\begin{proof}
Que $G\in AC$ est l'exemple~\ref{ex:integrale-ac} ; donc $G\in BV$ et $G'$ existe \pp{}
(corollaire~\ref{cor:bv}). Il reste à établir $G'=g$ \pp. On prolonge $g$ par $0$ hors de $[a,b]$.

\emph{Étape 1 : $g$ bornée, $\abs g\le M$.} Alors $G$ est $M$-lipschitzienne
($\abs{G(x)-G(y)}\le M\abs{x-y}$), donc dérivable \pp{} avec $\abs{G'}\le M$. Posons
$\varphi_n(x)=n\bigl(G(x+\tfrac1n)-G(x)\bigr)=n\int_x^{x+1/n}g$. Alors $\abs{\varphi_n}\le M$ et
$\varphi_n\to G'$ \pp. Par convergence dominée, pour tout $c\in[a,b]$,
\[
\int_a^c G'=\lim_n\int_a^c\varphi_n
=\lim_n n\Bigl[\int_c^{c+1/n}\!\!G-\int_a^{a+1/n}\!\!G\Bigr]=G(c)-G(a),
\]
la dernière égalité venant de la continuité de $G$ (valeur moyenne de $G$ sur un intervalle de longueur
$1/n$ autour de $c$, resp. $a$). Donc $\int_a^c(G'-g)=\int_a^c G'-\bigl(G(c)-G(a)\bigr)=0$ pour tout
$c$, et le lemme~\ref{lem:nul} donne $G'=g$ \pp.

\emph{Étape 2 : $g\ge0$, $g\in L^1$.} Pour $M>0$, posons $g_M=\min(g,M)$ (bornée) et
$R_M(x)=\int_a^x(g-g_M)$. Comme $g-g_M\ge0$, $R_M$ est croissante, donc $R_M'\ge0$ \pp{}
(théorème~\ref{thm:mono}). Or $G=\int g=\int g_M+R_M$, d'où par l'étape 1
\[
G'=g_M+R_M'\ge g_M\quad\pp.
\]
Ceci valant pour tout $M$, en faisant $M\to\infty$ : $G'\ge g$ \pp. On intègre : $\int_a^b G'\ge
\int_a^b g=G(b)-G(a)$. Mais $G$ est croissante (car $g\ge0$), donc $\int_a^b G'\le G(b)-G(a)$
(théorème~\ref{thm:mono}). Les deux inégalités forcent $\int_a^b(G'-g)=0$ avec $G'-g\ge0$, donc
$G'=g$ \pp.

\emph{Étape 3 : cas général.} Écrire $g=g^+-g^-$ avec $g^\pm\in L^1$, $g^\pm\ge0$, et appliquer
l'étape 2 à chacun : $\bigl(\int_a^x g^\pm\bigr)'=g^\pm$ \pp, d'où par différence $G'=g$ \pp.
\end{proof}

% ============================================================
\section{Volet II : la formule de Newton--Leibniz}
% ============================================================

\begin{theoreme}[TFA de Lebesgue]\label{thm:volet2}
Pour $F:[a,b]\to\R$, les assertions suivantes sont équivalentes :
\begin{enumerate}[label=\textup{(\roman*)},leftmargin=2.4em,itemsep=1pt]
  \item $F\in AC[a,b]$ ;
  \item il existe $g\in L^1$ avec $F(x)=F(a)+\int_a^x g$ pour tout $x$ ;
  \item $F$ est dérivable \pp, $F'\in L^1$, et $F(x)-F(a)=\int_a^x F'$ pour tout $x\in[a,b]$.
\end{enumerate}
Dans ce cas $g=F'$ \pp.
\end{theoreme}

\begin{proof}
$(iii)\Rightarrow(ii)$ est immédiat (prendre $g=F'$). $(ii)\Rightarrow(i)$ est l'exemple~\ref{ex:integrale-ac}.

$(ii)\Rightarrow(iii)$ : par le théorème~\ref{thm:volet1}, $F'=\bigl(F(a)+\int_a^x g\bigr)'=g$ \pp,
donc $F'\in L^1$ et $F(x)-F(a)=\int_a^x g=\int_a^x F'$.

$(i)\Rightarrow(iii)$ : c'est le cœur. Comme $F\in AC\subset BV$ (proposition~\ref{prop:hier}), $F'$
existe \pp{} et $F'\in L^1$ (corollaire~\ref{cor:bv}). Posons $G(x)=\int_a^x F'$. Par le
théorème~\ref{thm:volet1}, $G\in AC$ et $G'=F'$ \pp. La différence $H=F-G$ est alors $AC$ (différence
de deux fonctions $AC$) et vérifie $H'=F'-G'=0$ \pp. Le lemme de rigidité~\ref{lem:rigidite} donne $H$
constante ; comme $H(a)=F(a)-0=F(a)$, on a $F(x)-G(x)=F(a)$, soit
\[
F(x)-F(a)=G(x)=\int_a^x F'\qquad\forall x.\qedhere
\]
\end{proof}

\keybox{Le message central}{%
La formule $\int_a^b F'=F(b)-F(a)$ vaut \emph{exactement} pour les fonctions absolument continues.
« Continue » ne suffit pas ; « $BV$ » ne suffit pas ; « dérivable \pp » ne suffit pas. L'escalier de
Cantor appartient à toutes ces classes intermédiaires \emph{sauf} $AC$, et c'est précisément pourquoi
il échappe au TFA.}

\subsection{La caractérisation de Banach--Zaretsky}

La proposition~\ref{prop:luzin} livre gratuitement une seconde preuve de la non-absolue-continuité de
Cantor, et éclaire ce qui lui manque.

\begin{corollaire}\label{cor:cantor-N}
L'escalier de Cantor n'est pas absolument continu, car il viole la propriété (N) : il envoie $K$
(négligeable) sur $[0,1]$ (de mesure $1$), tandis que toute fonction $AC$ préserve la négligeabilité
(proposition~\ref{prop:luzin}).
\end{corollaire}

Le tableau se complète par le théorème de Banach--Zaretsky, dont on démontre la direction utile et
dont on énonce la réciproque : \emph{$F$ est $AC$ si et seulement si $F$ est continue, à variation
bornée, et vérifie la propriété (N).} La direction directe résulte de la proposition~\ref{prop:hier}
($AC\Rightarrow$ continue, $BV$) et de la proposition~\ref{prop:luzin} ($AC\Rightarrow$ (N)). La
réciproque, seule à faire appel à un résultat hors de notre boîte à outils (la formule de l'indicatrice
de Banach $\int_\R \#F^{-1}(y)\,dy=\Var_a^b F$), assure qu'aucun autre contre-exemple structurel que
le type « Cantor » ne peut exister parmi les fonctions continues $BV$.

% ============================================================
\section{Décomposition d'une fonction à variation bornée}
% ============================================================

Les contre-exemples ne sont pas des accidents : ce sont les composantes d'une décomposition canonique.

\subsection{Construction directe}

\begin{theoreme}[décomposition en trois parties]\label{thm:decomp}
Toute fonction croissante $F:[a,b]\to\R$ s'écrit
\[
F=F_{\mathrm{ac}}+F_{\mathrm{s}}+F_{\mathrm{saut}},
\]
où $F_{\mathrm{ac}}$ est absolument continue croissante, $F_{\mathrm{saut}}$ est une fonction de sauts
croissante (croissance concentrée sur un ensemble dénombrable), et $F_{\mathrm{s}}$ est croissante,
\emph{continue}, de dérivée nulle \pp{} (partie \emph{singulière continue}). Par différence, tout
$F\in BV$ se décompose de même.
\end{theoreme}

\begin{proof}
\emph{Partie de sauts.} $F$ croissante n'a qu'un ensemble dénombrable $D=\{t_j\}$ de discontinuités,
avec sauts $\sigma_j=F(t_j^+)-F(t_j^-)\ge0$ et $\sum_j\sigma_j\le F(b)-F(a)<\infty$. Posons
\[
F_{\mathrm{saut}}(x)=\sum_{t_j<x}\sigma_j+\bigl(F(x)-F(x^-)\bigr)\mathbf 1_{x\in D},
\]
croissante, dont les sauts reproduisent exactement ceux de $F$. Alors $\Phi:=F-F_{\mathrm{saut}}$ est
croissante et \emph{continue} (ses sauts s'annulent).

\emph{Parties absolument continue et singulière.} $\Phi$ croissante est dérivable \pp{} avec
$\Phi'\ge0$, $\Phi'\in L^1$ (théorème~\ref{thm:mono}). Posons
$F_{\mathrm{ac}}(x)=\int_a^x\Phi'$ (absolument continue, croissante, exemple~\ref{ex:integrale-ac})
et $F_{\mathrm{s}}=\Phi-F_{\mathrm{ac}}$. Par le théorème~\ref{thm:volet1}, $F_{\mathrm{ac}}'=\Phi'$
\pp, donc $F_{\mathrm{s}}'=\Phi'-\Phi'=0$ \pp. De plus $F_{\mathrm{s}}$ est continue (différence de
deux fonctions continues) et croissante : en effet, pour $x<y$,
$F_{\mathrm s}(y)-F_{\mathrm s}(x)=\bigl(\Phi(y)-\Phi(x)\bigr)-\int_x^y\Phi'\ge0$ par le
théorème~\ref{thm:mono} appliqué à $\Phi$ sur $[x,y]$. C'est la partie singulière continue.
\end{proof}

\begin{center}
\renewcommand{\arraystretch}{1.35}
\begin{tabular}{@{}>{\raggedright}p{3.4cm} c c >{\raggedright\arraybackslash}p{4.6cm}@{}}
\toprule
\textbf{Composante} & \textbf{Continue ?} & $\displaystyle\int F'=\Delta F$ ? & \textbf{Prototype} \\
\midrule
absolument continue & oui & \textbf{oui} & $\int_a^x g,\ g\in L^1$ \\
singulière continue & oui & non ($F'=0$ \pp) & \textbf{escalier de Cantor} \\
de sauts & non & non & fonction en escalier \\
\bottomrule
\end{tabular}
\end{center}

\noindent
La partie singulière continue est exactement l'obstruction au TFA parmi les fonctions \emph{continues} :
tout ensemble de Cantor négligeable porte une mesure de probabilité continue singulière, d'où une
famille entière de contre-exemples du même type que l'escalier.

\subsection{Le point de vue des mesures : Lebesgue--Radon--Nikodym}\label{sub:lrn}

La construction précédente est « à la main ». Le \emph{pourquoi} conceptuel apparaît en passant aux
mesures. À une fonction croissante $F$ (prise continue à droite pour fixer les idées) on associe son
unique \emph{mesure de Lebesgue--Stieltjes} $\mu_F$, mesure borélienne finie sur $[a,b]$ caractérisée
par
\[
\mu_F(]x,y])=F(y)-F(x),\qquad a\le x\le y\le b.
\]
Décomposer $F$ revient à décomposer $\mu_F$. Deux dichotomies \emph{orthogonales} sont à l'œuvre :
l'une relative à $\lambda$ (absolue continuité \emph{vs} singularité), l'autre intrinsèque (présence
\emph{vs} absence d'atomes). Leur superposition produit exactement trois briques.

\paragraph{Première dichotomie : absolue continuité contre singularité.}
Pour deux mesures $\mu,\nu$ sur un même espace mesurable $(X,\mathcal M)$ :
\begin{itemize}[leftmargin=1.6em,itemsep=1pt]
  \item $\mu$ est \emph{absolument continue} par rapport à $\nu$, noté $\mu\ll\nu$, si
  $\nu(E)=0\Rightarrow\mu(E)=0$ ;
  \item $\mu$ et $\nu$ sont \emph{étrangères}, noté $\mu\perp\nu$, s'il existe une partition
  $X=A\sqcup B$ ($\mathcal M$-mesurable) avec $\mu(B)=0$ et $\nu(A)=0$ : chacune vit là où l'autre ne
  charge rien.
\end{itemize}

\begin{theoreme}[Lebesgue--Radon--Nikodym]\label{thm:lrn}
Soient $\mu$ et $\lambda$ deux mesures \emph{finies} sur $(X,\mathcal M)$.
\begin{enumerate}[label=\textup{(\roman*)},leftmargin=2.4em,itemsep=1pt]
  \item \textup{(Décomposition de Lebesgue.)} $\mu$ s'écrit de manière \emph{unique}
  $\mu=\mu_{\mathrm{ac}}+\mu_{\perp}$ avec $\mu_{\mathrm{ac}}\ll\lambda$ et $\mu_{\perp}\perp\lambda$.
  \item \textup{(Radon--Nikodym.)} Il existe une fonction mesurable $g\ge0$, unique à égalité
  $\lambda$-\pp{} près, telle que $\mu_{\mathrm{ac}}(E)=\int_E g\,d\lambda$ pour tout $E\in\mathcal M$.
  On l'appelle \emph{densité} et on note $g=\dfrac{d\mu_{\mathrm{ac}}}{d\lambda}$.
\end{enumerate}
\end{theoreme}

\begin{proof}[Démonstration (von Neumann)]
Posons la mesure finie $\varphi=\mu+\lambda$ et travaillons dans l'espace de Hilbert $L^2(\varphi)$.
La forme linéaire $T(f)=\int_X f\,d\mu$ y est continue : par Cauchy--Schwarz,
\[
\abs{T(f)}\le\int_X\abs f\,d\mu\le\int_X\abs f\,d\varphi\le\varphi(X)^{1/2}\,\norm{f}_{L^2(\varphi)}.
\]
Le théorème de représentation de Riesz--Fréchet fournit donc $h\in L^2(\varphi)$ avec
\begin{equation}\label{eq:rn1}
\int_X f\,d\mu=\int_X f\,h\,d\varphi\qquad\text{pour tout }f\in L^2(\varphi).
\end{equation}
En prenant $f=\mathbf 1_E$ : $\mu(E)=\int_E h\,d\varphi$ ; comme $0\le\mu(E)\le\varphi(E)$, on a
$0\le h\le1$ $\varphi$-\pp. Puisque $d\varphi=d\mu+d\lambda$, l'égalité \eqref{eq:rn1} se réécrit
\begin{equation}\label{eq:rn2}
\int_X f\,(1-h)\,d\mu=\int_X f\,h\,d\lambda\qquad\text{pour tout }f\in L^2(\varphi).
\end{equation}
Séparons $X=A\sqcup B$ avec $A=\{0\le h<1\}$ et $B=\{h=1\}$, et posons
$\mu_{\mathrm{ac}}(E)=\mu(E\cap A)$, $\mu_{\perp}(E)=\mu(E\cap B)$.

\emph{$\mu_{\perp}\perp\lambda$.} Dans \eqref{eq:rn2} avec $f=\mathbf 1_B$ : le membre de gauche est
$\int_B(1-h)\,d\mu=0$ (car $h=1$ sur $B$), donc $\int_B h\,d\lambda=\lambda(B)=0$. Ainsi $\lambda$ ne
charge pas $B$, tandis que $\mu_{\perp}$ est concentrée sur $B$ : elles sont étrangères.

\emph{$\mu_{\mathrm{ac}}\ll\lambda$, avec densité.} Appliquons \eqref{eq:rn2} à
$f=(1+h+\cdots+h^n)\mathbf 1_E$. À gauche, $\int_E(1-h^{n+1})\,d\mu$ ; à droite,
$\int_E(h+h^2+\cdots+h^{n+1})\,d\lambda$. Quand $n\to\infty$, sur $A$ on a $h<1$ donc
$1-h^{n+1}\uparrow1$, et sur $B$ le terme de gauche est nul : par convergence monotone, le membre de
gauche tend vers $\mu(E\cap A)=\mu_{\mathrm{ac}}(E)$. À droite, la série croît vers
$g:=\sum_{k\ge1}h^k=\dfrac{h}{1-h}$ sur $A$ (et l'on pose $g=0$ sur $B$, licite car $\lambda(B)=0$) ;
par convergence monotone à nouveau,
\[
\mu_{\mathrm{ac}}(E)=\int_E g\,d\lambda.
\]
Donc $\mu_{\mathrm{ac}}\ll\lambda$ et $g\ge0$ est sa densité ; $g\in L^1(\lambda)$ car
$\mu_{\mathrm{ac}}(X)<\infty$.

\emph{Unicité.} Si $\mu=\mu_1+\nu_1=\mu_2+\nu_2$ avec $\mu_i\ll\lambda$ et $\nu_i\perp\lambda$, alors
$\eta:=\mu_1-\mu_2=\nu_2-\nu_1$ est à la fois $\ll\lambda$ et $\perp\lambda$. Étant $\perp\lambda$,
$\eta$ est concentrée sur un ensemble $\lambda$-négligeable $N$ ; étant $\ll\lambda$, elle s'annule sur
tout négligeable, donc $\eta(E)=\eta(E\cap N)=0$ pour tout $E$ : $\eta=0$. L'unicité de $g$ modulo
$\lambda$-\pp{} en découle.
\end{proof}

\begin{remarque}
Le cas $\sigma$-fini s'en déduit par exhaustion ($X=\bigsqcup_n X_n$, $\mu,\lambda$ finies sur chaque
$X_n$), et le cas signé par la décomposition de Jordan de la mesure. Enfin, appliqué à $\mu_F\ll\lambda$
avec $g=d\mu_{\mathrm{ac}}/d\lambda$, on a $g=F'$ $\lambda$-\pp{} : c'est le théorème de différentiation
de Lebesgue pour les mesures (les parties singulières ont une dérivée de Radon--Nikodym nulle \pp,
renforcement du théorème~\ref{thm:volet1} via un recouvrement de Vitali--Besicovitch). Pour la partie
absolument continue $\mu_{\mathrm{ac}}=g\,\lambda$, le volet~I (théorème~\ref{thm:volet1}) donne bien
$\bigl(\int_a^x g\bigr)'=g$ \pp.
\end{remarque}

\paragraph{Seconde dichotomie : atomes contre partie diffuse.}
Un point $x$ est un \emph{atome} de $\nu$ si $\nu(\{x\})>0$. La mesure $\nu$ est \emph{purement
atomique} si elle est portée par un ensemble dénombrable d'atomes, et \emph{diffuse} (ou
\emph{continue}) si $\nu(\{x\})=0$ pour tout $x$.

\begin{proposition}[décomposition atomique / diffuse]\label{prop:atom}
Toute mesure borélienne finie $\nu$ sur $\R$ s'écrit de manière unique
$\nu=\nu_{\mathrm{at}}+\nu_{\mathrm{cont}}$, avec $\nu_{\mathrm{at}}$ purement atomique et
$\nu_{\mathrm{cont}}$ diffuse.
\end{proposition}

\begin{proof}
L'ensemble des atomes $D=\{x:\nu(\{x\})>0\}$ est \emph{dénombrable} : pour chaque $m$, l'ensemble
$\{x:\nu(\{x\})>1/m\}$ a au plus $m\,\nu(\R)$ éléments (leurs masses, disjointes, se somment dans
$\nu(\R)$), et $D=\bigcup_m\{x:\nu(\{x\})>1/m\}$. Posons alors
$\nu_{\mathrm{at}}=\sum_{x\in D}\nu(\{x\})\,\delta_x=\nu(\,\cdot\cap D)$ et
$\nu_{\mathrm{cont}}=\nu(\,\cdot\setminus D)$. Pour tout $y$, $\nu_{\mathrm{cont}}(\{y\})=0$ (si
$y\in D$, on a retiré sa masse ; si $y\notin D$, $\nu(\{y\})=0$) : $\nu_{\mathrm{cont}}$ est diffuse.
L'unicité est immédiate car les atomes sont intrinsèques à $\nu$.
\end{proof}

\paragraph{Superposition : les trois briques.}
Appliquons Lebesgue--Radon--Nikodym à $\mu_F$, puis scindons la partie singulière $\mu_\perp$ selon la
proposition~\ref{prop:atom}, en $\mu_\perp=\mu_{\mathrm{at}}+\mu_{\mathrm{sc}}$. Deux observations
verrouillent l'articulation :
\begin{itemize}[leftmargin=1.6em,itemsep=1pt]
  \item \textbf{tout atome est singulier} : un singleton est $\lambda$-négligeable, donc
  $\mu_{\mathrm{at}}\perp\lambda$ ; la partie atomique loge \emph{entièrement} dans le singulier ;
  \item \textbf{la partie absolument continue est diffuse} : $\mu_{\mathrm{ac}}\ll\lambda$ et
  $\lambda(\{x\})=0$ entraînent $\mu_{\mathrm{ac}}(\{x\})=0$ ; elle n'a aucun atome.
\end{itemize}
Il n'y a donc pas de quatrième brique, et l'on obtient la décomposition canonique
\[
\mu_F=\underbrace{g\,\lambda}_{\ll\lambda\ (\text{diffuse})}
\;+\;\underbrace{\mu_{\mathrm{sc}}}_{\perp\lambda,\ \text{diffuse}}
\;+\;\underbrace{\sum_{j}\sigma_j\,\delta_{t_j}}_{\text{atomique}\ (\perp\lambda)},
\qquad g=F'\ \lambda\text{-\pp}.
\]

\begin{center}
\begin{tikzpicture}[
  every node/.style={font=\footnotesize},
  box/.style={draw=NavyBlue, rounded corners, align=center, inner sep=4pt, fill=BoxBlue},
  leaf/.style={draw=ForestGreen!65!black, rounded corners, align=center, inner sep=4pt},
  lab/.style={midway, font=\scriptsize\itshape, fill=white, inner sep=1pt}
]
  \node[box] (root) at (7,0) {$\mu_F$ \ (mesure de Stieltjes de $F$)};
  \node[box] (ac)   at (2.6,-2) {$\mu_{\mathrm{ac}}\ll\lambda$};
  \node[box] (si)   at (10.6,-2) {$\mu_{\perp}\perp\lambda$ \ (singulière)};
  \node[leaf] (acL) at (2.6,-4.4) {\textbf{absolument continue}\\ $\mu_{\mathrm{ac}}=F'\lambda$\\[1pt]
        $\to F_{\mathrm{ac}}(x)=\int_a^x F'$};
  \node[leaf] (at)  at (7.4,-4.4) {\textbf{atomique (sauts)}\\ $\sum_j\sigma_j\delta_{t_j}$\\[1pt]
        $\to F_{\mathrm{saut}}$};
  \node[leaf] (sc)  at (11.9,-4.4) {\textbf{singulière continue}\\ $\mu_{\mathrm{sc}}$ (type Cantor)\\[1pt]
        $\to F_{\mathrm{s}}$};
  \draw[gray!60] (root) -- (ac) node[lab,pos=0.55]{$\ll\lambda$};
  \draw[gray!60] (root) -- (si) node[lab,pos=0.55]{$\perp\lambda$};
  \draw[gray!60] (ac)   -- (acL);
  \draw[gray!60] (si)   -- (at) node[lab,pos=0.55]{atomique};
  \draw[gray!60] (si)   -- (sc) node[lab,pos=0.55]{diffuse};
\end{tikzpicture}
\end{center}

\noindent
\textbf{La réponse à « comment la partie singulière se scinde ».} La singularité (relation à $\lambda$)
et l'atomicité (propriété intrinsèque) sont indépendantes. La partie singulière $\mu_\perp$ rassemble
\emph{tout} ce qui est étranger à $\lambda$ ; on la trie alors selon la seconde dichotomie : sa
composante atomique $\mu_{\mathrm{at}}=\sum_j\sigma_j\delta_{t_j}$ (les \emph{sauts} de $F$, aux points
de discontinuité $t_j$, de masses $\sigma_j=F(t_j^+)-F(t_j^-)$) et sa composante diffuse
$\mu_{\mathrm{sc}}$, sans atome mais toujours étrangère à $\lambda$ (la partie \emph{singulière
continue}, portée par un négligeable, dont l'escalier de Cantor est le prototype via $\mu_c$). La
partie absolument continue, elle, est automatiquement diffuse : c'est pourquoi le monde « continu » se
partage encore en $\mathrm{ac}+\mathrm{sc}$, tandis que le monde « atomique » est tout entier singulier.

\paragraph{Retour aux fonctions (le dictionnaire).}
En intégrant chaque brique sur $]a,x]$, on retrouve terme à terme la décomposition du
théorème~\ref{thm:decomp} :
\[
F(x)-F(a)=\mu_F(]a,x])
=\underbrace{\int_a^x F'\,d\lambda}_{F_{\mathrm{ac}}(x)-F_{\mathrm{ac}}(a)}
\;+\;\underbrace{\mu_{\mathrm{sc}}(]a,x])}_{F_{\mathrm{s}}(x)-F_{\mathrm{s}}(a)}
\;+\;\underbrace{\sum_{a<t_j\le x}\sigma_j}_{F_{\mathrm{saut}}(x)-F_{\mathrm{saut}}(a)}.
\]
Le volet~II du TFA (théorème~\ref{thm:volet2}) se relit alors ainsi : $\int_a^x F'=F(x)-F(a)$ pour
\emph{tout} $x$ équivaut à la nullité des \emph{deux} briques singulières, c'est-à-dire à
$\mu_F=g\,\lambda$, soit $\mu_F\ll\lambda$ --- exactement l'absolue continuité de $F$.

\begin{center}
\renewcommand{\arraystretch}{1.35}
\begin{tabular}{@{}>{\raggedright}p{3.5cm} >{\raggedright}p{4.1cm} >{\raggedright\arraybackslash}p{6.4cm}@{}}
\toprule
\textbf{Brique de $\mu_F$} & \textbf{Nature} & \textbf{Fonction \& obstruction au TFA} \\
\midrule
$g\,\lambda$ & $\ll\lambda$, diffuse & $F_{\mathrm{ac}}=\int_a^x F'$ : aucune obstruction \\
$\mu_{\mathrm{sc}}$ & $\perp\lambda$, diffuse & $F_{\mathrm{s}}$ (Cantor) : $F'=0$ \pp, saut « lissé » \\
$\sum_j\sigma_j\delta_{t_j}$ & $\perp\lambda$, atomique & $F_{\mathrm{saut}}$ : croissance par sauts discrets \\
\bottomrule
\end{tabular}
\end{center}

% ============================================================
\section{Au-delà de Lebesgue : l'intégrale de Henstock--Kurzweil}
% ============================================================

Le volet II exige $F'\in L^1$. Or une fonction peut être dérivable \emph{partout} sans dérivée
Lebesgue-intégrable.

\begin{contreexemple}[dérivée partout, mais non $L^1$]\label{ce:hk}
Soit $F(x)=x^2\sin(1/x^2)$ pour $x\neq0$, $F(0)=0$. Alors $F$ est dérivable partout : pour $x\neq0$,
\[
F'(x)=2x\sin(1/x^2)-\tfrac{2}{x}\cos(1/x^2),
\]
et $F'(0)=\lim_{x\to0}\tfrac{F(x)}{x}=\lim_{x\to0}x\sin(1/x^2)=0$. Le terme $2x\sin(1/x^2)$ est borné
par $2\abs x$, donc intégrable. Pour l'autre, le changement de variable $t=1/x^2$ (soit
$x=t^{-1/2}$, $dx=-\tfrac12 t^{-3/2}\dt$) donne
\[
\int_0^1\abs*{\tfrac{2}{x}\cos(1/x^2)}\dx=\int_1^{\infty}2\,t^{1/2}\abs{\cos t}\cdot\tfrac12 t^{-3/2}\dt
=\int_1^{\infty}\frac{\abs{\cos t}}{t}\dt=+\infty,
\]
car $\abs{\cos t}$ est de moyenne $2/\pi>0$ et $\int_1^\infty \tfrac{dt}{t}=+\infty$. Donc
$\int_0^1\abs{F'}=+\infty$ : $F'\notin L^1$, et le volet~II ne s'applique pas, bien que $F$ soit
dérivable partout.
\end{contreexemple}

L'intégrale de Henstock--Kurzweil (dite \emph{de jauge}) répare cela. Une \emph{jauge} sur $[a,b]$ est
une fonction $\delta:[a,b]\to(0,\infty)$. Une subdivision pointée $P=\{(t_i,[x_{i-1},x_i])\}_{i=1}^n$
(avec $t_i\in[x_{i-1},x_i]$) est \emph{$\delta$-fine} si $[x_{i-1},x_i]\subset(t_i-\delta(t_i),t_i+\delta(t_i))$
pour tout $i$. On dit que $f$ est HK-intégrable d'intégrale $\mathcal A$ si, pour tout $\eps>0$, il
existe une jauge $\delta$ telle que toute subdivision pointée $\delta$-fine vérifie
$\abs{\sum_i f(t_i)(x_i-x_{i-1})-\mathcal A}<\eps$.

\begin{lemme}[Cousin]\label{lem:cousin}
Pour toute jauge $\delta$ sur $[a,b]$, il existe une subdivision pointée $\delta$-fine.
\end{lemme}

\begin{proof}
Supposons qu'il n'en existe pas. Coupons $[a,b]$ en deux moitiés : l'une au moins n'admet pas de
subdivision pointée $\delta$-fine (sinon on concaténerait). On construit ainsi une suite d'intervalles
emboîtés $[a_k,b_k]$, de longueur $(b-a)2^{-k}\to0$, aucun n'admettant de subdivision $\delta$-fine.
Soit $c$ le point commun $\bigcap_k[a_k,b_k]$. Pour $k$ assez grand,
$[a_k,b_k]\subset(c-\delta(c),c+\delta(c))$ ; alors la subdivision pointée à une seule cellule
$\{(c,[a_k,b_k])\}$ est $\delta$-fine --- contradiction.
\end{proof}

\begin{lemme}[« straddle »]\label{lem:straddle}
Soit $F$ dérivable en $t$ et $\eps>0$. Il existe $\eta>0$ tel que, pour tout intervalle $[u,v]$ avec
$u\le t\le v$ et $[u,v]\subset(t-\eta,t+\eta)$,
\[
\abs{F(v)-F(u)-F'(t)(v-u)}\le\eps(v-u).
\]
\end{lemme}

\begin{proof}
La dérivabilité en $t$ donne $\eta>0$ tel que $\abs{F(s)-F(t)-F'(t)(s-t)}\le\eps\abs{s-t}$ pour
$\abs{s-t}<\eta$. Appliqué à $s=v$ puis $s=u$ et ajouté (avec $t-u\ge0$, $v-t\ge0$) :
\[
\abs{F(v)-F(u)-F'(t)(v-u)}\le\eps(v-t)+\eps(t-u)=\eps(v-u).\qedhere
\]
\end{proof}

\begin{theoreme}[TFA de Henstock--Kurzweil]\label{thm:hk}
Si $F:[a,b]\to\R$ est dérivable en \emph{tout} point, alors $F'$ est HK-intégrable et
$\int_a^b F'=F(b)-F(a)$.
\end{theoreme}

\begin{proof}
Soit $\eps>0$. Pour chaque $t\in[a,b]$, le lemme~\ref{lem:straddle} fournit $\eta_t>0$ ; posons
$\delta(t)=\eta_t$, une jauge. Soit $P=\{(t_i,[x_{i-1},x_i])\}$ une subdivision pointée $\delta$-fine
(elle existe par le lemme de Cousin~\ref{lem:cousin}). Comme $t_i\in[x_{i-1},x_i]\subset
(t_i-\delta(t_i),t_i+\delta(t_i))$, le lemme~\ref{lem:straddle} donne
$\abs{F(x_i)-F(x_{i-1})-F'(t_i)(x_i-x_{i-1})}\le\eps(x_i-x_{i-1})$. En sommant, et comme les
$F(x_i)-F(x_{i-1})$ télescopent en $F(b)-F(a)$,
\[
\abs*{\sum_i F'(t_i)(x_i-x_{i-1})-\bigl(F(b)-F(a)\bigr)}
\le\sum_i\eps(x_i-x_{i-1})=\eps(b-a).
\]
C'est exactement la définition : $F'$ est HK-intégrable, d'intégrale $F(b)-F(a)$.
\end{proof}

C'est la forme la plus aboutie du volet Newton--Leibniz : \emph{aucune} hypothèse d'intégrabilité,
seule la dérivabilité en tout point. L'intégrale HK contient celle de Lebesgue (elles coïncident sur
$L^1$), tout en intégrant des dérivées non absolument intégrables comme celle du
contre-exemple~\ref{ce:hk}.

% ============================================================
\section{Synthèse}
% ============================================================

\begin{center}
\renewcommand{\arraystretch}{1.35}
\begin{tabular}{@{}>{\raggedright}p{5.0cm} >{\raggedright}p{4.3cm} c@{}}
\toprule
\textbf{Hypothèse sur $F$} & \textbf{Cadre} & $\int F'=\Delta F$ ? \\
\midrule
$F'$ continue & Riemann & \textbf{oui} \\
$F$ absolument continue & Lebesgue & \textbf{oui} \\
$F$ continue, $BV$ (Cantor) & Lebesgue & \textbf{non} \\
$F$ dérivable partout, $F'\in L^1$ & Lebesgue & \textbf{oui} \\
$F$ dérivable partout ($x^2\sin\frac1{x^2}$) & Lebesgue & \textbf{non} \\
$F$ dérivable partout & Henstock--Kurzweil & \textbf{oui} \\
\bottomrule
\end{tabular}
\end{center}

\medskip
\noindent
À chaque élargissement de l'intégrale (Riemann $\to$ Lebesgue $\to$ Henstock--Kurzweil) correspond un
affaiblissement des hypothèses rendant le TFA valide ; et à chaque étape, un contre-exemple précis ---
Volterra, l'escalier de Cantor, $x^2\sin(1/x^2)$ --- marque la frontière franchie et motive la
suivante. La condition charnière, celle qui gouverne le volet Newton--Leibniz au sens de Lebesgue,
est la \emph{continuité absolue}, exactement l'absence de partie singulière (continue ou de sauts) dans
la mesure $d\mu_F=F'\,d\lambda+\cdots$ associée à $F$.

\vspace{0.4cm}
\hrule
\vspace{0.3cm}
\noindent\textbf{\color{NavyBlue}Repères bibliographiques.}\ \
W.~Rudin, \emph{Real and Complex Analysis} (chap.~7) ;\ \
G.~B.~Folland, \emph{Real Analysis} (chap.~3) ;\ \
E.~M.~Stein \& R.~Shakarchi, \emph{Real Analysis} (chap.~3) ;\ \
F.~Riesz \& B.~Sz.-Nagy, \emph{Le\c{c}ons d'analyse fonctionnelle} ;\ \
R.~Gordon, \emph{The Integrals of Lebesgue, Denjoy, Perron, and Henstock}.

\end{document}
